% =============================================================================
%  main.tex
% =============================================================================
\documentclass{amr}

\usepackage{amsmath}

% Korrekturen für Textumbruch
\let\oldunderscore\_
\renewcommand{\_}{\oldunderscore\allowbreak}
\setlength{\emergencystretch}{1.5em}
\tolerance=1000

\begin{document}

% Titelseite & Verzeichnis
\maketitlepage
\tableofcontents
\clearpage

% Zusammenfassung
\abstractpage{
Die vorliegende Arbeit dokumentiert die Konzeption, den Entwurf und die experimentelle Validierung eines kostengünstigen autonomen mobilen Roboters (AMR) für strukturierte Innenräume.

Im Zentrum steht eine funktional geschichtete Architektur, die eine strikte Trennung zwischen zeitkritischen Regelungsaufgaben und rechenintensiven Navigationsfunktionen realisiert. Der fahrkritische Kern basiert auf einer verteilten Architektur mit zwei ESP32-S3-Mikrocontrollern, während die übergeordnete Lokalisierung und Kartierung sowie die Missionsplanung auf einem Raspberry Pi 5 unter ROS 2 ausgeführt werden. Ein wesentliches Merkmal des Systems ist die integrierte Sicherheits- und Freigabelogik, die den Vorrang des sicheren Zustands vor Navigationskommandos technisch sicherstellt. Die Validierung belegt eine reproduzierbare Grundbewegung sowie eine belastbare Hindernisvermeidung und schafft damit eine stabile Grundlage für künftige Erweiterungen in der intelligenten Interaktion.
}
\clearpage

% Hauptteil
\include{kap1}
\chapter{Grundlagen und Stand der Technik}

Dieses Kapitel klärt die technische Leitfrage der Arbeit: Welche Grundlagen tragen einen autonomen mobilen Roboter für den Transport von Kleinladungsträgern in einer veränderlichen Innenraumumgebung, und wie lassen sich diese Grundlagen in eine belastbare Systemarchitektur überführen? Die Darstellung folgt der Terminologie der Roadmap. Im Mittelpunkt stehen der Fahrkern, die Sensor- und Sicherheitsbasis, die verteilte ROS-2-Architektur, die Verfahren für Lokalisierung und Kartierung, die Navigation sowie die Interaktionsschichten für Bedienung und Sprache. Damit entsteht ein fachlicher Rahmen für die späteren Kapitel zur Anforderungsanalyse, zum Systemkonzept, zur Implementierung und zur Validierung.

\section{Autonome mobile Roboter in der Intralogistik}

Autonome mobile Roboter übernehmen Transportaufgaben in Umgebungen, deren Wege, Belegungen und Hindernisse nicht dauerhaft fest vorgegeben sind. Diese Eigenschaft unterscheidet sie von fahrerlosen Transportsystemen, die typischerweise an Leitlinien, Magnetbänder oder andere externe Führungsstrukturen gebunden sind. Für den Transport von Kleinladungsträgern in Werkstatt-, Labor- oder Wohnungsumgebungen ist diese Unterscheidung grundlegend, weil die Nutzbarkeit nicht nur von der Fahrfunktion, sondern vor allem von der Anpassungsfähigkeit an wechselnde Randbedingungen abhängt.

Der fachliche Kern eines AMR liegt deshalb nicht in einem einzelnen Sensor oder in einer einzelnen Softwarebibliothek, sondern in der Kopplung mehrerer Funktionen: Bewegungserzeugung, Zustandsschätzung, Umgebungswahrnehmung, Pfadplanung, Bahnverfolgung und Sicherheitsreaktion. Die Roadmap ordnet diese Funktionen dem Fahrkern, der Sensor- und Sicherheitsbasis, der Lokalisierung und Kartierung sowie der Navigation zu. Bedienung, Telemetrie und Sprachsteuerung bilden davon getrennte Interaktionsschichten. Diese Trennung reduziert Komplexität und verhindert, dass Bedienfunktionen unkontrolliert in sicherheits- oder fahrkritische Abläufe eingreifen.

Referenzplattformen aus Forschung und Lehre zeigen regelmäßig eine hierarchische Architektur. Mikrocontroller übernehmen hardwarenahe und zeitkritische Aufgaben, während leistungsfähigere Rechner Kartierung, Lokalisierung und Navigation ausführen. Die vorliegende Arbeit folgt diesem Grundmuster, verschärft jedoch die funktionale Trennung: Der Fahrkern und die Sensor- und Sicherheitsbasis laufen auf getrennten ESP32-S3-Knoten, während der Raspberry Pi 5 die hostseitigen ROS-2-Funktionen trägt. Diese Aufteilung dient nicht primär der Funktionsvielfalt, sondern der Entkopplung zeitkritischer Abläufe von blockierenden Sensor- oder Kommunikationszugriffen.

\section{Fahrkern: Differentialantrieb, Kinematik und Odometrie}

Der Fahrkern umfasst Antrieb, Odometrie und Grundbewegung. Für die vorliegende Plattform ist ein Differentialantrieb relevant. Zwei unabhängig angetriebene Räder erzeugen Translation und Rotation, indem sich ihre Winkelgeschwindigkeiten gezielt unterscheiden. Das Modell ist einfach genug für eine echtzeitfähige Regelung und zugleich hinreichend aussagekräftig für die Anbindung an Lokalisierung und Navigation.

Die Vorwärtskinematik beschreibt, wie aus den gemessenen Radwinkelgeschwindigkeiten $\omega_R$ und $\omega_L$ die translatorische Geschwindigkeit $v$ und die Giergeschwindigkeit $\omega$ des Roboters entstehen. Mit dem Radradius $r$ und der Spurbreite $b$ gilt:

$$v = \frac{r}{2} \left(\omega_R + \omega_L\right)$$

$$\omega = \frac{r}{b} \left(\omega_R - \omega_L\right)$$

Die Inverskinematik bildet den umgekehrten Weg ab. Sie wandelt eine Sollvorgabe aus dem Navigationssystem in Sollwerte für linkes und rechtes Rad um:

$$\omega_L = \frac{v - \omega \cdot \frac{b}{2}}{r}$$

$$\omega_R = \frac{v + \omega \cdot \frac{b}{2}}{r}$$

Diese Gleichungen koppeln den hostseitigen Navigationsstack direkt an die hardwarenahe Antriebsregelung. Damit aus dieser Kopplung reproduzierbare Bewegung entsteht, muss der Fahrkern systematische Fehler begrenzen. Besonders relevant sind Abweichungen im effektiven Raddurchmesser, in der Spurbreite, im Schlupfverhalten und in Totzonen der Motoransteuerung. Die Roadmap formuliert deshalb quantitative Kriterien für Geradeausfahrt, Rotation, Stoppverhalten und Wiederholbarkeit. Erst ein charakterisierter Fahrkern kann als belastbare Grundlage für die nachfolgenden Ebenen dienen.

Die Odometrie integriert Radbewegungen über die Zeit und liefert damit eine lokale Schätzung der Eigenbewegung. Ihr Vorteil liegt in der hohen Aktualisierungsrate und der unmittelbaren Verfügbarkeit. Ihre Schwäche liegt in der Fehlerakkumulation. Systematische Anteile, etwa das Verhältnis der effektiven Raddurchmesser $E_d$ und die effektive Spurbreite $E_b$, wirken über längere Fahrstrecken direkt auf die Positions- und Winkelschätzung. Die UMBmark-Kalibrierung adressiert genau diese Fehlerquellen durch definierte quadratische Fahrmanöver in beiden Drehrichtungen. Das Verfahren reduziert systematische Odometriefehler deutlich und schafft damit eine bessere Ausgangsbasis für Sensorfusion und Kartierung.

\section{Sensor- und Sicherheitsbasis}

Die Sensor- und Sicherheitsbasis umfasst alle Signale, die den Fahrkern stabilisieren, Gefährdungen früh erkennen oder den Energiezustand überwachen. In der Roadmap zählen dazu insbesondere IMU, Ultraschallsensor, Cliff-Sensor und Batterieüberwachung. Diese Ebene ergänzt den Fahrkern um Messgrößen, die aus den Radencodern allein nicht hervorgehen. Gleichzeitig bildet sie die erste technische Schutzschicht gegen unsichere Zustände.

Die inertiale Messeinheit erfasst Drehraten und Beschleunigungen. Für mobile Robotik ist vor allem die Drehrate um die Hochachse relevant, weil sie kurzfristige Orientierungsänderungen auch dann abbildet, wenn Radschlupf die encoderbasierte Schätzung verfälscht. Dem steht eine typische Schwäche inertialer Sensoren gegenüber: Bias, Drift und temperaturabhängige Abweichungen verschieben die Schätzung über die Zeit. Deshalb genügt eine IMU nicht als alleinige Quelle der Pose. Sie gewinnt ihren Wert erst in der Fusion mit Odometrie und gegebenenfalls weiteren absoluten Referenzen.

Der Front-Ultraschallsensor bildet den Nahbereich ab und ergänzt die flächige Umgebungswahrnehmung um eine einfache Distanzinformation im direkten Vorfeld. Die Batterieüberwachung erfasst Spannung und Strom, damit Unterspannung früh erkannt und Schutzreaktionen ausgelöst werden können. Beide Funktionen sind für einen reproduzierbaren Betrieb wichtiger als ihre vergleichsweise einfache Sensorik vermuten lässt: Ohne Energieüberwachung verliert jede Bewegungsbewertung an Vergleichbarkeit, und ohne Nahbereichserfassung steigt das Risiko, dass kurzreichweitige Hindernisse oder Betriebsgrenzen zu spät erkannt werden.

Eine besondere Rolle übernimmt die hardwarenahe Kantenerkennung. Der Cliff-Sensor adressiert eine Gefährdung, bei der die Reaktionszeit entscheidend ist. Rein softwarebasierte Navigationsalgorithmen arbeiten mit Verarbeitungsketten, Scheduling und Kommunikationswegen, die zwar im Normalbetrieb ausreichen, aber keinen unmittelbaren Schutz gegen abrupte Kanten garantieren. Deshalb ordnet die Architektur das Signal des Cliff-Sensors einer Sicherheitslogik zu, die Bewegungsfreigaben sofort sperren und eine Notbremsung auslösen kann. Im Systemkontext handelt es sich nicht um Komfortfunktion, sondern um eine übergeordnete Schutzmaßnahme.

\section{Verteilte Software-Architektur mit ROS 2 und micro-ROS}

ROS 2 bildet die Middleware der hostseitigen Systemebene. Das Modell basiert auf ROS-2-Knoten, Topics, Diensten, Aktionen und Launch-Dateien. Der fachliche Nutzen dieser Struktur liegt in der klaren Trennung von Funktionen. Ein Knoten verarbeitet genau abgegrenzte Aufgaben, veröffentlicht Zustände oder Messwerte über Topics und lässt sich mit anderen Knoten über standardisierte Schnittstellen koppeln. Dadurch bleibt das Gesamtsystem erweiterbar, obwohl Wahrnehmung, Lokalisierung, Navigation, Bedienung und Audio parallel laufen.

micro-ROS erweitert diesen Ansatz auf Mikrocontroller. Die beiden ESP32-S3-Knoten treten damit nicht als isolierte Hilfscontroller auf, sondern als eingebundene Teile der Gesamtarchitektur. Für die Arbeit ist diese Eigenschaft wesentlich, weil der Fahrkern und die Sensor- und Sicherheitsbasis Zustände direkt in das ROS-2-System einspeisen, ohne dass eine proprietäre Nebenarchitektur entstehen muss. Die Kommunikation erfolgt über serielle UART-Verbindungen beziehungsweise USB-CDC. Gegenüber drahtlosen Verbindungen bietet diese Kette reproduzierbarere Laufzeiten und eine geringere Störanfälligkeit. Für einen Regelkreis mit $50\,\mathrm{Hz}$ ist diese Entscheidung keine Nebenbedingung, sondern ein Stabilitätsfaktor.

Die verteilte Aufteilung folgt einem klaren Architekturprinzip. Der Drive-Knoten verarbeitet Antriebsregelung und encodernahe Funktionen. Der Sensor-Knoten verarbeitet I2C-basierte Sensorik, Batteriedaten und sicherheitsnahe Signale. Der Raspberry Pi 5 übernimmt Kartierung, Lokalisierung, Navigation, Benutzeroberfläche und höhere Interaktionsfunktionen. Damit trennt die Architektur zeitkritische Abläufe von rechen- und kommunikationsintensiven Hostaufgaben. Gleichzeitig verhindert die physische Aufteilung, dass langsame Sensorzugriffe die Motorregelung blockieren. Diese Entkopplung ist ein direkter Beitrag zur Robustheit des Fahrkerns.

\section{Lokalisierung und Kartierung}

Lokalisierung und Kartierung beantworten zwei eng gekoppelte Fragen: Wo befindet sich der Roboter, und wie sieht die Umgebung aus? Für ein AMR mit veränderlicher Einsatzumgebung genügt weder reine Odometrie noch eine einmalig eingemessene Infrastruktur. Erforderlich ist eine Karten- und Transformationskette, die Sensordaten räumlich einordnet, Drift begrenzt und Wiederanlauf nach Neustart ermöglicht. Die Roadmap ordnet diese Aufgaben der Ebene „Lokalisierung und Kartierung" zu. Zu dieser Ebene gehören LiDAR, TF-Baum, SLAM-Verfahren, Kartenrepräsentation und Re-Lokalisierung.

Der LiDAR liefert Distanzmessungen in einer Ebene und erzeugt daraus einen Scan der Umgebung. Die geometrische Nutzbarkeit dieses Scans hängt nicht nur vom Sensor selbst ab, sondern auch von den zugehörigen Koordinatensystemen. Der TF-Baum verknüpft Basisrahmen, Sensorräume und Weltbezug. Fehler in diesen Transformationen erzeugen keine kleinen Schönheitsfehler, sondern systematische Kartenverzerrungen, instabile Re-Lokalisierung oder Sprünge zwischen Bezugssystemen. Deshalb gehört die Extrinsik des Sensors fachlich zur Kartierungsgrundlage und nicht erst zur späteren Optimierung.

Für die simultane Lokalisierung und Kartierung kommt ein SLAM-Verfahren zum Einsatz. Die SLAM Toolbox kombiniert Scan-Matching mit einer graphbasierten Optimierung und erzeugt daraus eine konsistente Karte, obwohl Odometrie und Einzelmessungen fehlerbehaftet sind. Der fachliche Wert eines solchen Verfahrens liegt nicht nur in der Kartenerzeugung, sondern in der Korrektur lokaler Schätzfehler durch globale Konsistenzbedingungen. Das Ergebnis ist eine wiederverwendbare Karte, auf deren Basis spätere Zielanfahrten überhaupt erst reproduzierbar werden.

\section{Navigation, Sicherheitslogik und Missionslogik}

Navigation beginnt erst dann sinnvoll, wenn Fahrkern sowie Lokalisierung und Kartierung hinreichend belastbar arbeiten. Fachlich umfasst die Navigation die globale Pfadplanung, die lokale Bahnverfolgung, die Hindernisberücksichtigung und das Recovery-Verhalten. In der Roadmap ist diese Ebene explizit von der Bedienung getrennt. Diese Trennung verhindert, dass einzelne Bedienereignisse oder Sensoreffekte ungefiltert in Fahrbefehle münden.

Nav2 stellt dafür einen modularen Navigationsstack bereit. Ein globaler Planer erzeugt einen Weg durch die Karte. Ein lokaler Regler verfolgt diesen Weg unter Berücksichtigung der Fahrzeugkinematik und der aktuellen Hindernissituation. Costmaps bilden statische und dynamische Hindernisse als Kostenfelder ab. Behavior Trees koordinieren diese Funktionsblöcke und verknüpfen Standardverhalten mit Fehlerreaktionen wie Anhalten, Rücksetzen oder Neuversuch. Gegenüber starren Zustandsmaschinen bietet diese Struktur eine bessere Erweiterbarkeit bei gleichzeitig klaren Entscheidungswegen.

Für die vorliegende Arbeit reicht reine Navigation jedoch nicht aus. Zwischen Interaktion und Bewegungsfreigabe ist eine Freigabelogik erforderlich. Sie bewertet, ob ein Befehl zulässig ist, in welchem Betriebsmodus sich das System befindet und ob Sicherheitsbedingungen erfüllt sind. Erst danach entsteht ein Missionskommando, etwa eine Zielanfahrt oder ein definierter Betriebswechsel. Die eigentliche Bewegungsausgabe bleibt weiterhin an Navigation, Fahrkern und Sicherheitslogik gebunden. Damit gilt als Architekturregel:

$$\text{Interaktion} \rightarrow \text{Freigabelogik} \rightarrow \text{Missionskommando} \rightarrow \text{Navigation} \rightarrow \text{Fahrkern}$$

Diese Kette begrenzt das Risiko unsicherer Direktansteuerung und macht Entscheidungswege nachvollziehbar. Sie ist damit nicht nur ein Softwaremuster, sondern ein Sicherheitsprinzip.

\section{Bedien- und Leitstandsebene}

Die Bedien- und Leitstandsebene stellt Beobachtbarkeit und Eingriffsmöglichkeiten bereit, ohne die innere Fahrlogik aufzulösen. In der Roadmap gehören dazu Dashboard, Telemetrie, manuelle Kommandos, Audio-Rückmeldungen und weitere Betriebsanzeigen. Fachlich entspricht diese Ebene einem Leitstandkonzept: Zustände werden sichtbar, Diagnoseinformationen werden gesammelt, Bedienhandlungen werden protokollierbar und Betriebsarten werden nachvollziehbar umgeschaltet.

Für ein MINT-orientiertes Entwicklungsprojekt hat diese Ebene einen doppelten Nutzen. Erstens verbessert sie den Betrieb, weil Akkuzustand, Sensordaten, Kamerabilder, Fahrzustände und Fehlermeldungen zentral beobachtbar werden. Zweitens verbessert sie die Validierung, weil Versuche reproduzierbarer dokumentiert und Fehlersituationen klarer zugeordnet werden können. Eine Benutzeroberfläche ist damit kein Zusatzmodul am Rand, sondern ein Werkzeug zur technischen Beherrschung des Gesamtsystems. Die Roadmap nennt dafür bereits konkrete Schnittstellen wie \texttt{/cmd\_vel}, \texttt{/servo\_cmd}, \texttt{/hardware\_cmd} und \texttt{/audio/play}. Im Architekturkontext gilt jedoch auch hier: Bedienung darf Fahrfunktionen auslösen, aber sie darf die Sicherheitslogik nicht umgehen.

\section{Sprachschnittstelle als sichere Interaktionsschicht}

Die Sprachschnittstelle erweitert die Bedien- und Leitstandsebene um natürliche Eingaben. Sie ersetzt weder Navigation noch Sicherheitslogik. Ihr fachlicher Zweck besteht darin, Sprachsignale in klar definierte Befehlsabsichten zu überführen und nur freigegebene Operationen an die übrige Systemarchitektur weiterzugeben. Die Roadmap ordnet dafür das ReSpeaker Mic Array v2.0, Sprach-zu-Text-Verarbeitung, Intent-Erkennung, Befehlsmultiplexing und Sprachrückmeldung einer eigenen Teilarchitektur zu.

Das Sicherheitsprinzip dieser Ebene lautet: Kein Sprachbefehl darf direkt in eine rohe Motoransteuerung übersetzt werden. Zulässig ist nur die Kette aus Sprachbefehl, Intent, Freigabelogik und Missionskommando. Sichere Sofortkommandos wie „Stopp" oder „Halt" dürfen ausschließlich einen sicheren Halt auslösen. Betriebsmodus-Kommandos dürfen Zustände wechseln, aber keine direkten Geschwindigkeiten setzen. Missionskommandos dürfen Zielpunkte oder Aktionen anfordern, nicht jedoch die Schutzlogik außer Kraft setzen. Informationskommandos dienen der Zustandsabfrage und unterstützen die Bedienung bei geringem Risiko. Damit bleibt die Sprachschnittstelle eine Interaktionsschicht und wird nicht zu einem unsicheren Parallelpfad in den Fahrkern.

\section{Zusammenfassung}

Der Stand der Technik zeigt kein einzelnes Schlüsselmodul, sondern eine funktional geschichtete Architektur. Der Fahrkern erzeugt reproduzierbare Grundbewegung. Die Sensor- und Sicherheitsbasis ergänzt den Fahrkern um robuste Mess- und Schutzfunktionen. ROS 2 und micro-ROS koppeln Host- und Mikrocontroller-Ebene in einer verteilten Struktur. Lokalisierung und Kartierung schaffen den räumlichen Bezug. Navigation überführt Kartenwissen in sichere Zielanfahrten. Bedien- und Leitstandsebene sowie Sprachschnittstelle erweitern das System um beobachtbare und kontrollierte Interaktion, ohne die Sicherheitslogik zu unterlaufen. Aus dieser Schichtung folgt die zentrale Konsequenz für die weitere Arbeit: Zusätzliche Funktionen erhöhen nur dann den Systemwert, wenn der Fahrkern, die Sicherheitslogik und die Freigabelogik zuvor belastbar geschlossen sind.

\chapter{Anforderungsanalyse}

\section{Leitfrage, Einsatzszenario und Systemgrenze}

Die Leitfrage dieses Kapitels lautet: Welche Anforderungen muss ein kostengünstiges autonomes mobiles Robotersystem erfüllen, damit Fahrkern, Sensor- und Sicherheitsbasis, Lokalisierung und Kartierung, Navigation sowie Bedien- und Leitstandsebene in einem Innenraum belastbar zusammenarbeiten?

Das Einsatzszenario dieser Arbeit ist der autonome Transport kleiner Nutzlasten zwischen definierten Zielpunkten in einem strukturierten Innenraum. Die Referenzumgebung entspricht einer wohnungs- oder laborähnlichen Fläche mit Engstellen, Möbeln, Wänden und wechselnden Hindernissen. Damit verlagert sich der Schwerpunkt gegenüber einem klassischen Industrie-Szenario: Nicht maximaler Materialdurchsatz, sondern reproduzierbare Navigation, sichere Bewegungsfreigabe und nachvollziehbare Systemreaktionen stehen im Vordergrund.

Der Anwendungsfall gliedert sich in fünf aufeinanderfolgende Schritte. Zunächst erzeugt das System eine Karte oder nutzt eine vorhandene Karte zur Re-Lokalisierung. Anschließend nimmt die Bedien- und Leitstandsebene ein Missionskommando entgegen, etwa eine Zielanfahrt zu einem definierten Punkt. Danach plant die Navigation einen globalen Pfad und übergibt Sollgrößen an den Fahrkern. Während der Fahrt überwacht die Sensor- und Sicherheitsbasis den Nahbereich, den Energiezustand und sicherheitsnahe Ereignisse. Nach Erreichen des Zielpunkts meldet die Benutzeroberfläche das Ergebnis und erlaubt einen Folgeauftrag oder einen sicheren Halt.

Für die Anforderungsableitung ist die Systemgrenze entscheidend. Bestandteil der Kernarchitektur sind Fahrkern, Sensor- und Sicherheitsbasis, Lokalisierung und Kartierung, Navigation, Sicherheitslogik, Freigabelogik sowie die Bedien- und Leitstandsebene. Eine Sprachschnittstelle gehört zur geplanten Ausbauarchitektur. Sie darf nur freigegebene Missionskommandos erzeugen und keine rohe Motoransteuerung auslösen. Nicht Gegenstand der Kernvalidierung sind Außenbetrieb, Treppenfahrt, freies semantisches Weltmodell oder vollautomatisches Greifen.

\section{Technische Randbedingungen und Restriktionen}

\subsection{Hardware- und Plattformrandbedingungen}

Die Systemarchitektur folgt einer verteilten Aufteilung zwischen Host-Ebene und Mikrocontroller-Ebene. Der Raspberry Pi 5 übernimmt hostseitige Funktionen wie Lokalisierung und Kartierung, Navigation, Benutzeroberfläche, Telemetrie, Audio-Rückmeldung und spätere Sprachverarbeitung. Zwei XIAO-ESP32-S3-Mikrocontroller übernehmen hardwarenahe Aufgaben. Der Drive-Knoten bildet den Fahrkern mit Motoransteuerung, Encoder-Auswertung und Odometrie. Der Sensor-Knoten bildet die Sensor- und Sicherheitsbasis mit IMU, Ultraschall, Kanten-Erkennung und Batterieüberwachung. Diese Rollenverteilung entspricht der Roadmap-Terminologie und trennt zeitkritische Regelung von blockierenden Sensorzugriffen.

Als Sensorik dienen ein LiDAR für die 2D-Umgebungserfassung, eine inertiale Messeinheit für Drehraten- und Beschleunigungsdaten, ein Ultraschallsensor für den Nahbereich, ein Cliff-Sensor für Kanten und eine Batterieüberwachung für den Energiezustand. Die Referenz-Roadmap ordnet diese Komponenten explizit den Phasen Fahrkern, Sensor- und Sicherheitsbasis sowie Lokalisierung und Kartierung zu.

\subsection{Software- und Kommunikationsrandbedingungen}

Die Host-Software basiert auf ROS 2 Humble. Da das Zielsystem Debian Trixie nutzt, läuft der ROS-2-Stack in einer Container-Umgebung. Die Mikrocontroller werden über micro-ROS als ROS-2-Teilnehmer eingebunden. Für die Kommunikation zwischen Host und Mikrocontrollern kommt eine serielle UART-Verbindung über USB-CDC zum Einsatz. Die Anforderung folgt nicht aus Bequemlichkeit, sondern aus dem Bedarf an reproduzierbaren Laufzeiten. Eine drahtlose Verbindung würde die zeitliche Streuung der Kommandokette erhöhen und die Bewertung des Fahrkerns erschweren. Die bisherige Fassung von Kapitel 3 setzt für den Regelkreis einen Arbeitstakt von $50\,\mathrm{Hz}$ an.

Die Softwarearchitektur trennt Interaktion, Freigabelogik, Missionskommando, Navigation und Fahrkern. Daraus folgt eine verbindliche Befehlskette:

$$\text{Interaktion} \rightarrow \text{Freigabelogik} \rightarrow \text{Missionskommando} \rightarrow \text{Navigation} \rightarrow \text{Fahrkern}$$

Jede Anforderung, die Kommandos erzeugt oder verarbeitet, muss diese Kette einhalten.

\subsection{Umgebungs- und Sicherheitsrandbedingungen}

Die Arbeit betrachtet einen ebenen Innenraum mit bekannten oder kartierbaren Strukturen. Zulässig sind Wände, Möbel, Türen, Engstellen und bewegliche Hindernisse wie Personen. Nicht zulässig sind Außenflächen, nasse Untergründe, starke Steigungen oder ungesicherte Treppenfahrten. Eine erkannte Kante ist kein zu passierendes Hindernis, sondern ein Abbruchkriterium für die Bewegung.

Die Zielgeschwindigkeit beträgt $0{,}4\,\mathrm{m/s}$. Dieser Wert begrenzt die maximale Dynamik und reduziert zugleich die Anforderungen an Sensorreichweite, Bremsweg und Kollisionsvermeidung. Für Zielanfahrten gelten als Orientierungswerte eine laterale Zielabweichung von höchstens $0{,}10\,\mathrm{m}$ und ein Orientierungsfehler von höchstens $8^\circ$. Diese Größen dienen im weiteren Verlauf als Referenz für funktionale Anforderungen und Akzeptanzkriterien.

\section{Funktionale Anforderungen}

\subsection{F01 — Fahrkern mit reproduzierbarer Grundbewegung}

Der Fahrkern muss lineare und rotatorische Bewegung reproduzierbar ausführen. Dazu gehören Motoransteuerung, Encoder-Auswertung, Odometrie und die Umsetzung von Geschwindigkeitskommandos. Der Fahrkern muss eine Geradeausfahrt über $1\,\mathrm{m}$ mit kleinem Seitenfehler, eine Rotation um $360^\circ$ mit reproduzierbarem Winkelfehler und einen sicheren Stopp ohne ungewolltes Nachlaufen unterstützen. Diese Forderung leitet sich direkt aus Phase 1 der Roadmap ab.

\subsection{F02 — Sensor- und Sicherheitsbasis mit priorisierten Schutzfunktionen}

Die Sensor- und Sicherheitsbasis muss IMU-Daten, Nahbereichsdaten, Batteriedaten und Kanteninformationen in definierter Form bereitstellen. Sicherheitsnahe Signale müssen gegenüber Komfort- oder Diagnosefunktionen priorisiert werden. Eine erkannte Kante muss die Bewegungsfreigabe sperren und einen sicheren Halt auslösen. Die Batterieüberwachung muss einen kritischen Energiezustand melden, bevor die Fahrfunktion unkontrolliert ausfällt. Diese Forderung leitet sich aus Phase 2 der Roadmap ab.

\subsection{F03 — Lokalisierung und Kartierung für den Innenraum}

Das System muss aus LiDAR-Daten, Odometrie und Transformationsbeziehungen eine konsistente Karte erzeugen und sich nach einem Neustart in dieser Karte wieder lokalisieren können. Der TF-Baum muss widerspruchsfrei bleiben. Die Kartenauflösung beträgt als Zielwert $5\,\mathrm{cm}$. Wiederholte Kartierungen desselben Raums dürfen keine ausgeprägten Doppelkonturen erzeugen. Diese Forderung leitet sich aus Phase 3 der Roadmap ab.

\subsection{F04 — Navigation mit Missionslogik und Recovery-Verhalten}

Die Navigation muss definierte Zielpunkte sicher anfahren. Dazu gehören globale Pfadplanung, lokale Bahnverfolgung, Hindernisberücksichtigung und nachvollziehbares Recovery-Verhalten. Die Navigation darf nur freigegebene Missionskommandos ausführen. Einzelne Benutzereingaben oder experimentelle Zusatzmodule dürfen keine direkte Umgehung der Sicherheitslogik erzeugen. Für die Kernvalidierung gilt als Zielwert die erfolgreiche Durchführung von $10$ definierten Zielanfahrten ohne Kollision. Diese Forderung leitet sich aus Phase 4 der Roadmap ab.

\subsection{F05 — Bedien- und Leitstandsebene als Betriebswerkzeug}

Die Bedien- und Leitstandsebene muss Zustände sichtbar machen und definierte Eingriffe erlauben. Dazu gehören Telemetrie, Statusanzeigen, Kamerabild oder Videostream, manuelle Kommandos in freigegebenen Betriebsarten und Audio-Rückmeldungen. Die Benutzeroberfläche dient nicht nur der Bedienung, sondern auch der Diagnose und Versuchsdokumentation. Die vorhandenen Schnittstellen \texttt{/cmd\_vel}, \texttt{/servo\_cmd}, \texttt{/hardware\_cmd} und \texttt{/audio/play} bilden dafür die technische Grundlage. Diese Forderung leitet sich aus Ebene B und Phase 5 der Roadmap ab.

\subsection{F06 — Freigabelogik mit sicheren Zustandsübergängen}

Die Freigabelogik muss festlegen, welche Kommandos im jeweiligen Betriebszustand zulässig sind. Sie muss Kommandos freigeben, blockieren oder in sichere Reaktionen umsetzen. Ein Stopp-Kommando hat Vorrang vor allen Fahr- und Missionskommandos. Bedienkommandos dürfen Navigation und Fahrkern beeinflussen, aber nicht die Schutzmechanismen umgehen. Diese Forderung ergibt sich aus der Terminologie-Norm der Roadmap mit den Begriffen Sicherheitslogik, Freigabelogik und Missionskommando.

\subsection{F07 — Sprachschnittstelle als anschlussfähige Erweiterung}

Die Sprachschnittstelle soll Sprachsignale in definierte Intents überführen und daraus freigegebene Missionskommandos ableiten. Die Sprachschnittstelle darf keine rohen Geschwindigkeitsbefehle direkt an den Fahrkern senden. Sichere Sofortkommandos wie „Stopp" dürfen ausschließlich einen sicheren Halt anfordern. Diese Forderung gehört nicht zur Kernvalidierung, aber zur thematischen Erweiterung der Projektarbeit. Die Roadmap ordnet das ReSpeaker Mic Array v2.0, Sprach-zu-Text-Verarbeitung, Intent-Erkennung und Text-zu-Sprache der Ebene der intelligenten Interaktion zu.

\section{Nichtfunktionale Anforderungen}

\subsection{N01 — Deterministische Verarbeitung im Fahrkern}

Der Fahrkern muss Sollwerte und Messwerte mit hinreichend konstanter Zykluszeit verarbeiten. Für die Motorregelung gilt ein Arbeitstakt von $50\,\mathrm{Hz}$. Die zeitliche Streuung darf den Regelkreis nicht instabil machen. Blockierende I2C-Zugriffe oder nicht deterministische Nebenaufgaben dürfen den Fahrkern nicht unterbrechen. Die bisherige Ausgangsfassung fordert hierfür einen Jitter von weniger als $2\,\mathrm{ms}$.

\subsection{N02 — Modulare und wartbare Systemstruktur}

Die Architektur muss Funktionen in getrennten ROS-2-Knoten, Topics und Launch-Dateien abbilden. Fahrkern, Sensor- und Sicherheitsbasis, Lokalisierung und Kartierung, Navigation sowie Bedien- und Leitstandsebene müssen fachlich getrennt bleiben. Diese Modularität erleichtert Fehlersuche, Erweiterung und Wiederverwendung.

\subsection{N03 — Robuste Kommunikation und Fehlerbehandlung}

Die serielle micro-ROS-Kommunikation muss Verbindungsabbrüche erkennen und einen definierten Wiederanlauf unterstützen. Kommunikationsfehler dürfen nicht zu unkontrollierter Weiterfahrt führen. Stattdessen muss das System in einen sicheren Zustand wechseln oder auf einen sicheren Halt zurückfallen. Die Ausgangsfassung von Kapitel 3 nennt den automatischen Wiederanlauf der UART-Kommunikation ausdrücklich als nichtfunktionale Anforderung.

\subsection{N04 — Beobachtbarkeit und Nachvollziehbarkeit}

Zustände, Sensordaten, Betriebsarten und Fehlerereignisse müssen so bereitgestellt werden, dass Versuche reproduzierbar ausgewertet werden können. Dazu gehören Telemetrie, Protokollierung und eine klar strukturierte Benutzeroberfläche. Eine Bewertung ohne sichtbare Mess- und Zustandsbasis ist nicht belastbar.

\subsection{N05 — Erweiterbarkeit der Interaktionsschicht}

Die Architektur muss spätere Erweiterungen wie Sprachschnittstelle, Audio-Rückmeldung oder zusätzliche semantische Module aufnehmen können, ohne die Kernlogik des Fahrkerns zu verändern. Erweiterbarkeit bedeutet in diesem Kontext nicht beliebige Offenheit, sondern kontrollierte Kopplung über Freigabelogik und Missionskommandos.

\section{Priorisierte Anforderungsliste nach MoSCoW}

\begin{table}[htbp]
  \centering
  \caption{Priorisierte Anforderungsliste nach MoSCoW}
  \small
  \begin{tabularx}{\textwidth}{l L c c L}
    \toprule
    \textbf{ID} & \textbf{Anforderung} & \textbf{Typ} & \textbf{Prio.} & \textbf{Akzeptanzkriterium} \\
    \midrule
    F01 & Fahrkern mit reproduzierbarer Grundbewegung & F & M & Geradeausfahrt über $1\,\mathrm{m}$ reproduzierbar; Rotation um $360^\circ$ reproduzierbar; kein ungewolltes Nachlaufen nach Stopp. \\
    \addlinespace
    F02 & Sensor- und Sicherheitsbasis mit priorisierten Schutzfunktionen & F & M & IMU, Ultraschall, Kanten-Erkennung und Batterieüberwachung liefern nutzbare Signale; erkannte Kante führt reproduzierbar zum sicheren Halt. \\
    \addlinespace
    F03 & Lokalisierung und Kartierung für den Innenraum & F & M & Karte mit Zielauflösung von $5\,\mathrm{cm}$ erzeugbar; Re-Lokalisierung nach Neustart möglich; keine ausgeprägten Doppelkonturen. \\
    \addlinespace
    F04 & Navigation mit Missionslogik und Recovery-Verhalten & F & M & $10$ definierte Zielanfahrten ohne Kollision; Zielradius dokumentiert; Fehlfahrten und Recovery-Verhalten nachvollziehbar protokolliert. \\
    \addlinespace
    F05 & Bedien- und Leitstandsebene als Betriebswerkzeug & F & M & Telemetrie, Statusanzeige, manuelle Kommandos und Audio-Rückmeldung sind verfügbar; Zustände sind über die Benutzeroberfläche nachvollziehbar. \\
    \addlinespace
    F06 & Freigabelogik mit sicheren Zustandsübergängen & F & M & Nicht freigegebene Kommandos werden blockiert; Stopp hat Vorrang; Kommandokette folgt dem Schema Interaktion $\rightarrow$ Freigabelogik $\rightarrow$ Missionskommando. \\
    \addlinespace
    F07 & Sprachschnittstelle als anschlussfähige Erweiterung & F & C & Sprachbefehl wird in Intent und Missionskommando überführt; direkte rohe Motoransteuerung aus Sprache ist ausgeschlossen. \\
    \addlinespace
    N01 & Deterministische Verarbeitung im Fahrkern & NF & M & Regelzyklus mit $50\,\mathrm{Hz}$; Jitter kleiner als $2\,\mathrm{ms}$; Sensorzugriffe blockieren die Motorregelung nicht dauerhaft. \\
    \addlinespace
    N02 & Modulare und wartbare Systemstruktur & NF & S & Funktionen sind in getrennten ROS-2-Knoten und Launch-Dateien organisiert; fachliche Ebenen bleiben entkoppelt. \\
    \addlinespace
    N03 & Robuste Kommunikation und Fehlerbehandlung & NF & M & Kommunikationsausfall wird erkannt; System fällt in sicheren Zustand; Wiederanlauf der UART-basierten micro-ROS-Kette ist möglich. \\
    \addlinespace
    N04 & Beobachtbarkeit und Nachvollziehbarkeit & NF & S & Zustände, Sensordaten und Fehlerereignisse sind protokollierbar und über die Benutzeroberfläche zugänglich. \\
    \addlinespace
    N05 & Erweiterbarkeit der Interaktionsschicht & NF & S & Audio- und Sprachfunktionen lassen sich ergänzen, ohne Fahrkern oder Sicherheitslogik strukturell umzubauen. \\
    \bottomrule
  \end{tabularx}
\end{table}

Die MoSCoW-Priorisierung folgt einer klaren Regel. Must-Anforderungen bilden die Kernvalidierung der Arbeit. Should-Anforderungen verbessern Wartbarkeit, Diagnose und spätere Systemqualität. Could-Anforderungen markieren geplante Erweiterungen, insbesondere an der Interaktionsschicht. Damit entsteht eine Anforderungsbasis, die die Roadmap-Themen in eine projektarbeitstaugliche Form überführt und zugleich die spätere Validierung in Kapitel 6 vorbereitet.

\chapter{Systemkonzept und Entwurf}

Dieses Kapitel beantwortet die Entwurfsfrage der Arbeit: Wie lässt sich ein Innenraum-AMR für Kartierung, Zielanfahrt und sichere Bedienung so strukturieren, dass der Fahrkern reproduzierbar arbeitet, sicherheitsrelevante Signale Vorrang erhalten und spätere Erweiterungen wie Vision oder Sprachschnittstelle anschlussfähig bleiben? Die Darstellung überführt die Anforderungen aus Kapitel 3 in eine konkrete Zielarchitektur. Im Mittelpunkt stehen die Konzeptauswahl, die funktionale Partitionierung, die Daten- und Kommandokette sowie die Einordnung von Lokalisierung und Kartierung, Navigation, Bedien- und Leitstandsebene und intelligenter Interaktion.

\section{Entwurfsziele und Auswahlkriterien}

Das Systemkonzept folgt vier Entwurfszielen. Erstens muss der Fahrkern zeitlich stabil arbeiten, damit Geschwindigkeitsvorgaben reproduzierbar in Bewegung umgesetzt werden. Zweitens muss die Sensor- und Sicherheitsbasis sicherheitsnahe Signale priorisieren, damit eine erkannte Kante oder ein kritischer Betriebszustand eine sofortige Schutzreaktion auslösen kann. Drittens muss die Architektur Funktionen sauber trennen, damit Kartierung, Navigation, Telemetrie und Audio die hardwarenahe Regelung nicht stören. Viertens muss die Struktur spätere Erweiterungen aufnehmen, ohne die Kernlogik des Fahrkerns umzubauen.

Aus diesen Zielen ergeben sich die Auswahlkriterien für den Entwurf:

\begin{table}[htbp]
  \centering
  \caption{Entwurfsziele und deren Bedeutung}
  \small
  \begin{tabularx}{\textwidth}{l X}
    \toprule
    \textbf{Kriterium} & \textbf{Bedeutung für den Entwurf} \\
    \midrule
    Determinismus & Der Fahrkern muss Sollwerte und Messwerte mit konstanter Zykluszeit verarbeiten. \\
    \addlinespace
    Priorisierte Sicherheitsreaktion & Sicherheitslogik und Freigabelogik müssen Fahrkommandos überstimmen können. \\
    \addlinespace
    Modulare Struktur & Fahrkern, Sensor- und Sicherheitsbasis, Lokalisierung und Kartierung, Navigation sowie Bedienung müssen getrennt bleiben. \\
    \addlinespace
    Innenraumeignung & Kinematik, Sensorik und Navigation müssen für enge Räume, Kurven und variable Hindernisse geeignet sein. \\
    \addlinespace
    Erweiterbarkeit & Benutzeroberfläche, Audio, Vision und Sprachschnittstelle müssen anschließbar bleiben. \\
    \bottomrule
  \end{tabularx}
\end{table}

Diese Kriterien verschieben die Entwurfsentscheidung bewusst weg von maximaler Rechenleistung oder maximaler Funktionsvielfalt. Vorrang erhalten stattdessen klare Zuständigkeiten, stabile Datenpfade und nachvollziehbare Zustandsübergänge.

\section{Konzeptauswahl}

Der morphologische Kasten reduziert den Variantenraum auf die Teilfunktionen, die den Systemcharakter am stärksten prägen. Die ausgewählte Lösung erfüllt die Anforderungen des Innenraum-AMR mit dem geringsten Integrationsrisiko.

\begin{table}[htbp]
  \centering
  \caption{Morphologischer Kasten zur Konzeptauswahl}
  \small
  \begin{tabularx}{\textwidth}{l L L L}
    \toprule
    \textbf{Teilfunktion} & \textbf{Alternative A} & \textbf{Alternative B} & \textbf{Alternative C} \\
    \midrule
    Host-Recheneinheit & \textbf{Raspberry Pi 5 mit optionalem Hailo-8L-Beschleuniger} & Jetson Nano & Intel NUC \\
    \addlinespace
    Low-Level-Steuerung & \textbf{2 x ESP32-S3} & 1 x ESP32-S3 & STM32F767ZI \\
    \addlinespace
    Kommunikation Host-Low-Level & \textbf{micro-ROS über UART} & USB mit proprietärer Kopplung & WLAN über UDP \\
    \addlinespace
    Antriebskonzept & \textbf{Differentialantrieb} & Mecanum-Antrieb & Ackermann-Lenkung \\
    \addlinespace
    Lokalisierung und Kartierung & \textbf{SLAM Toolbox} & Cartographer & GMapping \\
    \addlinespace
    Lokaler Navigationsregler & \textbf{Regulated Pure Pursuit} & Dynamic Window Approach & MPPI \\
    \bottomrule
  \end{tabularx}
\end{table}

Die Wahl des Raspberry Pi 5 folgt der Anforderung, ROS 2, LiDAR, Kartierung, Navigation, Benutzeroberfläche und Audio auf einer einzigen Host-Plattform zusammenzuführen. Ein zusätzlicher Hailo-8L-Beschleuniger bleibt als Erweiterung für Vision-Aufgaben möglich, ist jedoch nicht Voraussetzung für den Kernbetrieb. Der Host bildet damit die Rechenbasis für Lokalisierung und Kartierung, Navigation sowie Bedien- und Leitstandsebene.

Die Aufteilung der Low-Level-Ebene auf zwei ESP32-S3 ist die zentrale Entwurfsentscheidung. Der Drive-Knoten übernimmt ausschließlich Fahrkern und Odometrie. Der Sensor-Knoten übernimmt I2C-basierte Sensorik, Batterieüberwachung und sicherheitsnahe Signale. Diese physische Trennung senkt das Risiko, dass blockierende Sensorzugriffe oder langsame Peripherieoperationen die Motorregelung stören.

Für die Kommunikation zwischen Host und Low-Level-Ebene wurde micro-ROS über UART gewählt. Die serielle Kopplung reduziert den Integrationsaufwand, erlaubt eine direkte Einbindung der Mikrocontroller in das ROS-2-System und vermeidet die zusätzliche Unsicherheit einer drahtlosen Übertragung im Regelpfad.

Zusätzlich verbindet ein CAN-Bus mit $1\,\mathrm{Mbit/s}$ (ISO 11898, SN65HVD230-Transceiver) die beiden ESP32-S3-Knoten direkt miteinander. Dieser Dual-Path erlaubt sicherheitsrelevante Signale wie Cliff-Erkennung und Unterspannungsabschaltung unabhängig vom Host und micro-ROS Agent zu übertragen. Der CAN-Bus bildet damit einen redundanten Sicherheitspfad neben der UART-basierten micro-ROS-Kommunikation.

Der Differentialantrieb ist für enge Innenräume die zweckmäßige Wahl. Das Konzept benötigt wenig mechanische Komplexität, unterstützt Rotation auf engem Raum und lässt sich mit überschaubarem Rechenaufwand modellieren. Mecanum- und Ackermann-Konzepte würden den mechanischen und regelungstechnischen Aufwand erhöhen, ohne für die Zielumgebung einen zwingenden Vorteil zu liefern.

Für Lokalisierung und Kartierung wurde die SLAM Toolbox gewählt, weil sie Kartenaufbau und Re-Lokalisierung in einer ROS-2-nahen Toolchain abbildet. Für die lokale Bahnverfolgung wurde Regulated Pure Pursuit gewählt, weil das Verfahren kinematisch einfache Plattformen gut unterstützt und Geschwindigkeiten in Kurven oder bei Hindernisannäherung begrenzen kann.

\section{Zielarchitektur des Gesamtsystems}

Das Zielbild folgt der Roadmap mit drei Ebenen. Die Ebenen trennen fahrkritische Funktionen, Betriebsfunktionen und intelligente Interaktion.

\begin{table}[htbp]
  \centering
  \caption{Ebenen der Zielarchitektur}
  \small
  \begin{tabularx}{\textwidth}{l L L}
    \toprule
    \textbf{Ebene} & \textbf{Hauptfunktionen} & \textbf{Zentrale Recheneinheit} \\
    \midrule
    Ebene A – Fahrkern sowie Sensor- und Sicherheitsbasis & Antrieb, Odometrie, IMU, Ultraschall, Cliff-Sensor, Batterieüberwachung, Servo-Steuerung, Sicherheitsreaktion & 2 x ESP32-S3 und Raspberry Pi 5 \\
    \addlinespace
    Ebene B – Bedien- und Leitstandsebene & Benutzeroberfläche, Telemetrie, manuelle Kommandos, Videostream, Audio-Rückmeldungen & Raspberry Pi 5 \\
    \addlinespace
    Ebene C – Intelligente Interaktion & Sprachschnittstelle, semantische Interpretation, Vision, spätere multimodale Bedienung & Raspberry Pi 5 mit optionalem Beschleuniger \\
    \bottomrule
  \end{tabularx}
\end{table}

Die physische Ausführung ordnet den Ebenen feste Komponenten zu. Der Raspberry Pi 5 trägt die hostseitigen ROS-2-Knoten für Lokalisierung und Kartierung, Navigation, Benutzeroberfläche, Audio und spätere Interaktionsfunktionen. Ein ESP32-S3 bildet den Drive-Knoten des Fahrkerns. Ein weiterer ESP32-S3 bildet den Sensor-Knoten der Sensor- und Sicherheitsbasis. LiDAR, Kamera und Audio-Hardware sind an den Host angebunden.

Die Architektur trennt nicht nur Hardware, sondern auch Verantwortlichkeiten. Der Fahrkern setzt freigegebene Fahrvorgaben um. Die Sensor- und Sicherheitsbasis liefert Zustände und Schutzsignale. Lokalisierung und Kartierung sowie Navigation berechnen Bewegungsziele auf Kartenbasis. Die Bedien- und Leitstandsebene beobachtet und bedient das System. Die Sprachschnittstelle bleibt eine Interaktionsschicht oberhalb der Freigabelogik.

\section{Mechanisches und elektronisches Konzept}

Das mechanische Konzept basiert auf einem Differentialantrieb mit zwei angetriebenen Rädern und einer frei laufenden Stützstruktur. Der Aufbau unterstützt enge Kurvenradien und Richtungswechsel auf engem Raum. Für den Prototyp sind ein Raddurchmesser von $65{,}67\,\mathrm{mm}$ und eine kalibrierte Spurbreite von $178\,\mathrm{mm}$ vorgesehen. Diese Größen gehen direkt in Kinematik, Odometrie und spätere Kalibrierung ein.

Die Antriebseinheit nutzt JGA25-370-Motoren mit Encoder-Rückführung. Die Leistungsstufe bildet ein Cytron-MDD3A-Motortreiber. Die PWM-Frequenz beträgt $20\,\mathrm{kHz}$, um die Ansteuerung oberhalb des gut hörbaren Bereichs zu betreiben und gleichzeitig eine fein genug aufgelöste Stellgröße bereitzustellen.

Die Sensorik folgt der funktionalen Trennung des Gesamtsystems. Encoder erfassen die Radbewegung des Fahrkerns. Eine MPU6050 liefert inertiale Messgrößen. Ein Front-Ultraschallsensor erfasst den Nahbereich. Ein Cliff-Sensor erkennt Kanten. Eine INA260 überwacht Spannung und Strom des Energiesystems. I2C-basierte Baugruppen einschließlich der Servo-Steuerung über einen PCA9685-PWM-Treiber liegen vollständig auf dem Sensor-Knoten, damit der Fahrkern nicht durch Sensorzugriffe belastet wird.

Zusätzliche Hardware für die Bedien- und Leitstandsebene und die intelligente Interaktion wird bewusst vom Kernpfad getrennt. Eine frontseitige Kamera unterstützt Videostream und spätere Vision-Funktionen. Das ReSpeaker Mic Array v2.0 dient als Audioeingang für die Sprachschnittstelle. Ein Lautsprecherpfad ermöglicht definierte Audio-Rückmeldungen. Diese Komponenten erweitern die Beobachtbarkeit und Interaktion, ohne die Grundfahrt direkt zu steuern.

\section{Software-Architektur und funktionale Partitionierung}

Die Software-Architektur kombiniert ROS 2 auf dem Host mit micro-ROS auf den Mikrocontrollern. Auf dem Raspberry Pi 5 laufen die ROS-2-Knoten für Lokalisierung und Kartierung, Navigation, Benutzeroberfläche, Audio und Schnittstellenlogik. Auf den ESP32-S3 läuft FreeRTOS als harte Partitionierung der Low-Level-Aufgaben. Damit entsteht eine verteilte Architektur mit klarer Zuständigkeit pro Knoten.

Der Drive-Knoten verarbeitet ausschließlich fahrkritische Funktionen. Dazu gehören die Umsetzung von Geschwindigkeitsvorgaben, die Encoder-Auswertung, die Odometrie und die Motorregelung mit einem Arbeitstakt von $50\,\mathrm{Hz}$. Die hostseitige Navigation liefert translatorische und rotatorische Sollgrößen, der Drive-Knoten setzt diese über das kinematische Modell in Radgrößen um.

Der Sensor-Knoten verarbeitet die Sensor- und Sicherheitsbasis. Er liest IMU, Ultraschall, Batterieüberwachung und Kanten-Erkennung aus und veröffentlicht die Daten als ROS-2-nahe Informationen über micro-ROS. Zusätzlich liefert er sicherheitsnahe Zustände an die hostseitige Sicherheitslogik.

Die hostseitige Ebene bündelt Kartenaufbau, Re-Lokalisierung, Zielplanung, Benutzeroberfläche und Audio. Diese Aufteilung erlaubt eine klare Trennung zwischen hardwarenaher Reaktion und rechenintensiver Verarbeitung. Zugleich vermeidet sie eine proprietäre Nebenarchitektur, weil beide Mikrocontroller als eingebundene ROS-2-Teilnehmer auftreten.

Die Kommandokette bildet ein zentrales Entwurfselement. Sie verhindert, dass Bedienung oder Sprache rohe Motorbefehle direkt in den Fahrkern einspeisen. Zulässig ist nur eine freigegebene Kette:

\begin{align*}
  \text{Interaktion} &\rightarrow \text{Freigabelogik} \rightarrow \text{Missionskommando} \\
                     &\rightarrow \text{Navigation oder manuelle Bedienung} \rightarrow \text{/cmd\_vel} \rightarrow \text{Fahrkern}
\end{align*}

Nicht zulässig ist eine direkte Kette der Form:

$$\text{Sprachbefehl} \rightarrow \text{direkte Motoransteuerung}$$

Diese Trennung gibt der Sicherheitslogik und der Freigabelogik einen festen Ort in der Architektur. Ein Stopp-Kommando hat Vorrang vor Fahr- und Missionskommandos. Nicht freigegebene Eingaben werden blockiert oder in einen sicheren Zustand überführt.

\section{Datenfluss, Sicherheitslogik und Freigabelogik}

Die Daten- und Befehlsflüsse folgen dem Prinzip „Zustand nach oben, Freigabe nach unten". Sensorwerte, Odometrie und Diagnoseinformationen laufen von den Mikrocontrollern zum Host. Missionskommandos, Zielzustände und freigegebene Fahrbefehle laufen vom Host zum Fahrkern.

Für die Navigation entstehen mindestens drei logisch getrennte Befehlsquellen: autonome Fahrbefehle der Navigation, manuelle Fahrbefehle der Benutzeroberfläche und Stopp- oder Schutzkommandos der Sicherheitslogik. Die Architektur führt diese Quellen nicht unkontrolliert zusammen, sondern über eine Freigabelogik mit eindeutigem Vorrang. Die Sicherheitslogik steht oberhalb der regulären Fahrvorgaben. Eine erkannte Kante darf daher eine aktive Zielanfahrt unmittelbar unterbrechen.

Der Cliff-Sicherheitsmultiplexer bildet diese Regel technisch ab. Er verarbeitet den Status des Cliff-Sensors und übersteuert bei Bedarf eingehende Fahrbefehle. Das Ergebnis ist kein konkurrierender Fahrkanal, sondern eine übergeordnete Schutzinstanz. Damit trennt das System bewusst zwischen Navigationsfehlern, die Recovery-Verhalten auslösen können, und Schutzereignissen, die unmittelbar zum Halt führen müssen.

Die Freigabelogik regelt zusätzlich die Übergänge zwischen Betriebsarten. Navigation darf nur freigegebene Missionskommandos ausführen. Manuelle Eingriffe der Benutzeroberfläche dürfen den Fahrzustand beeinflussen, aber keine Schutzmechanismen umgehen. Sprachkommandos dürfen nur freigegebene Missionskommandos erzeugen. Dadurch bleibt die gesamte Befehlskette auch bei späteren Erweiterungen nachvollziehbar.

\section{Entwurf von Fahrkern, Lokalisierung und Kartierung sowie Navigation}

Der Fahrkern arbeitet mit dem kinematischen Modell des Differentialantriebs. Die hostseitige Ebene erzeugt eine Translationsgeschwindigkeit $v$ und eine Drehrate $\omega$. Der Drive-Knoten berechnet daraus die Radsollwerte:

$$\omega_L = \frac{v - \omega \cdot \frac{b}{2}}{r}$$

$$\omega_R = \frac{v + \omega \cdot \frac{b}{2}}{r}$$

mit dem Radradius $r$ und der Spurbreite $b$. Die konkrete Parametrierung der Motorregelung gehört zur Implementierung und Validierung. Im Systemkonzept reicht die Festlegung, dass der Fahrkern freigegebene Sollgrößen deterministisch und reproduzierbar umsetzen muss.

Für Lokalisierung und Kartierung sind zwei Betriebsarten vorgesehen. Im Kartierungsbetrieb erzeugt die SLAM Toolbox aus LiDAR-Daten, Odometrie und Transformationsbeziehungen eine Karte des Innenraums. Im Navigationsbetrieb nutzt das System eine vorhandene Karte, lokalisiert sich darin erneut und fährt Zielpunkte auf Kartenbasis an. Diese Trennung reduziert Komplexität, weil Kartenaufbau und autonome Zielanfahrt nicht dauerhaft denselben Optimierungszustand teilen müssen.

Die Navigation kombiniert globale Pfadplanung mit lokaler Bahnverfolgung. Globale Planer erzeugen eine Route auf Basis der Karte. Der lokale Regler verfolgt diese Route unter Berücksichtigung der Fahrzeugkinematik. Recovery-Verhalten bleibt Teil der Navigation, etwa bei blockiertem Weg oder ungünstiger Pose. Recovery ersetzt jedoch keine Sicherheitslogik. Eine Kante, ein kritischer Energiezustand oder eine gesperrte Freigabe sind Abbruchkriterien und keine regulären Navigationsprobleme.

\section{Bedien- und Leitstandsebene sowie Sprachschnittstelle}

Die Bedien- und Leitstandsebene ist kein Nebenprodukt des Entwurfs, sondern ein eigenes Systemelement. Sie stellt Telemetrie, Zustandsanzeige, Videostream, manuelle Kommandos und Audio-Rückmeldungen bereit. Damit unterstützt sie Diagnose, Versuchsdurchführung und den sicheren Betrieb. Die Benutzeroberfläche dient folglich nicht nur der Interaktion, sondern auch der Beobachtbarkeit des Systems.

Die Sprachschnittstelle wird als Erweiterung der Ebene der intelligenten Interaktion eingeordnet. Das ReSpeaker Mic Array v2.0 liefert Audiodaten an den Host. Nachgelagerte ROS-2-Knoten können Sprache in Text, Text in einen Intent und den Intent in ein freigegebenes Missionskommando überführen. Die Sprachschnittstelle bleibt damit oberhalb von Freigabelogik und Missionskommando. Sie ergänzt die Benutzeroberfläche, ersetzt aber weder Navigation noch Sicherheitslogik.

Die daraus entstehende Kette lautet:

\begin{align*}
  \text{Sprachbefehl} &\rightarrow \text{Intent} \rightarrow \text{Freigabelogik} \\
                      &\rightarrow \text{Missionskommando} \\
                      &\rightarrow \text{Navigation, Leitstand oder Audio}
\end{align*}

Durch diese Einordnung lässt sich die Sprachschnittstelle in die Projektarbeit aufnehmen, ohne den Kernnachweis des Fahr- und Navigationssystems zu verwässern.

\section{Ergebnis des Systementwurfs}

Der Entwurf überführt die Anforderungen des Innenraum-AMR in eine klar gegliederte Zielarchitektur. Die Dual-Node-Struktur trennt Fahrkern und Sensor- und Sicherheitsbasis. Der Raspberry Pi 5 bündelt Lokalisierung und Kartierung, Navigation, Benutzeroberfläche und spätere Interaktionsfunktionen. Sicherheitslogik und Freigabelogik sichern die Kommandokette gegen unkontrollierte Eingriffe. Damit liegt ein Systemkonzept vor, das den Kernbetrieb der Projektarbeit trägt und zugleich Erweiterungen wie Audio, Vision und Sprachschnittstelle geordnet einbindet.

\chapter{Implementierung}

Dieses Kapitel beantwortet die Umsetzungsfrage der Arbeit: Wie wurde die in Kapitel 4 entworfene Architektur technisch realisiert, so dass der Fahrkern reproduzierbar arbeitet, sicherheitsrelevante Signale Vorrang erhalten und Bedien- und Leitstandsfunktionen den Kernbetrieb nicht stören? Die Darstellung folgt der Roadmap-Struktur und trennt deshalb konsequent zwischen Fahrkern, Sensor- und Sicherheitsbasis, hostseitiger ROS-2-Integration, Sicherheitslogik, Bedien- und Leitstandsebene sowie anschlussfähigen Erweiterungen für intelligente Interaktion.

\section{Hardwareaufbau und elektrische Inbetriebnahme}

Die Hardwareimplementierung verfolgt zwei Ziele. Erstens muss der Fahrkern mechanisch und elektrisch so aufgebaut sein, dass Odometrie und Antriebsregelung reproduzierbare Mess- und Stellgrößen erhalten. Zweitens muss die Sensor- und Sicherheitsbasis so angebunden sein, dass sicherheitsnahe Signale unabhängig von rechenintensiven Host-Funktionen erfasst werden können.

\subsection{Chassis, Antrieb und Grundgeometrie}

Die mechanische Basis bildet ein Differentialantrieb mit zwei koaxial angeordneten Getriebemotoren und einer frei nachlaufenden Stützstruktur. Die präzise Ausrichtung der Motorachsen ist für die Odometrie wesentlich, weil eine fehlerhafte Spurbreite systematische Bahnabweichungen verursacht. Die Antriebsräder wurden deshalb spielfrei auf den Motorwellen montiert. Der verwendete Radradius beträgt

$$r = 32{,}835\,\mathrm{mm}$$

Dieser Wert entspricht dem kalibrierten Raddurchmesser von $65{,}67\,\mathrm{mm}$, der aus Bodentests mit Maßbandvergleich hervorging (siehe Kapitel 6).

Die Stützstruktur wurde so positioniert, dass ein möglichst großer Anteil der Fahrzeugmasse auf den angetriebenen Rädern lastet. Diese Maßnahme reduziert den Einfluss nicht modellierter Roll- und Schlupfeffekte bei Richtungswechseln.

\subsection{Verdrahtung der Dual-Node-Architektur}

Die Low-Level-Ebene wurde mit zwei physisch getrennten XIAO-ESP32-S3-Knoten aufgebaut. Diese Trennung folgt direkt dem Entwurf aus Kapitel 4: Der Drive-Knoten verarbeitet ausschließlich fahrkritische Funktionen, während der Sensor-Knoten I2C-basierte Peripherie und sicherheitsnahe Signale übernimmt. Die Verdrahtung bildet diese Zuständigkeiten eindeutig ab.

Am Drive-Knoten sind die Hall-Encoder mit den Phasen A und B an die Pins D4 und D5 (Phase A) sowie D8 und D9 (Phase B) angeschlossen. Die PWM-Signale für den Cytron-MDD3A-Motortreiber liegen im Dual-PWM-Modus an den Pins D0 bis D3 an. Die Signalleitungen der Encoder wurden räumlich getrennt von den PWM-führenden Motorleitungen verlegt, um kapazitive Einkopplungen der Schaltflanken zu verringern. Die PWM-Frequenz beträgt

$$f_{\mathrm{PWM}} = 20\,\mathrm{kHz}$$

Am Sensor-Knoten bindet der I2C-Bus mit SDA auf D4 und SCL auf D5 die MPU6050, den INA260 und den PCA9685 an. Ein MH-B-Kanten-Sensor ist direkt an einen digitalen GPIO-Pin angeschlossen. Diese Aufteilung vermeidet, dass I2C-Zugriffe den Fahrkern blockieren.

\subsection{Host-nahe Peripherie}

Der Raspberry Pi 5 bindet die hostseitigen Komponenten für Lokalisierung und Kartierung, Navigation, Bedien- und Leitstandsebene sowie spätere Interaktion an. Der RPLIDAR A1 ist per USB angeschlossen und liefert die Distanzdaten für Kartenaufbau und Navigation. Eine frontseitig montierte Sony-IMX296-Kamera ist über CSI angebunden. Für hardwarebeschleunigte Bildverarbeitung wurde ein Hailo-8L-Beschleuniger über ein M.2-HAT integriert. Der Audiopfad nutzt einen MAX98357A-Verstärker als Ausgabestufe.

Diese Peripherie liegt bewusst außerhalb des Low-Level-Regelpfads. Damit bleiben Kamera, Vision und Audio anschlussfähig, ohne die Zykluszeit der Motorregelung direkt zu beeinflussen.

\subsection{Spannungsversorgung und Massekonzept}

Die Spannungsversorgung trennt Rechenlast, Logik und Leistungspfad. Der Raspberry Pi 5 wird über einen Step-Down-Regler mit

$$U = 5\,\mathrm{V}$$

versorgt. Die beiden ESP32-S3-Knoten beziehen ihre Versorgung über USB; die Boards erzeugen die Logikspannung von

$$U = 3{,}3\,\mathrm{V}$$

über interne Regler. Der Motortreiber schaltet die Akkuspannung direkt auf die Antriebe. Der Leistungsbereich liegt damit im Bereich von

$$U = 10{,}8\,\mathrm{V}\ \text{bis}\ 12{,}6\,\mathrm{V}$$

Ein gemeinsamer Sternpunkt für die Masse reduziert Ausgleichsströme und erschwert Masseschleifen zwischen Host, Mikrocontrollern und Leistungsstufe.

\section{Firmware von Fahrkern sowie Sensor- und Sicherheitsbasis}

Die Firmware wurde in zwei getrennten PlatformIO-Projekten umgesetzt. Jedes Projekt läuft auf einem eigenen ESP32-S3-Knoten und nutzt FreeRTOS zur funktionalen und zeitlichen Partitionierung. Diese Trennung überführt die Architekturentscheidung aus Kapitel 4 in lauffähige Firmware.

\subsection{Drive-Knoten für Fahrkern und Odometrie}

Der Drive-Knoten realisiert den Fahrkern. Er verarbeitet Encoder-Signale, berechnet Odometrie und setzt freigegebene Geschwindigkeitsvorgaben in PWM-Signale für den Motortreiber um. Die Software ist in hardwarenahe und regelungstechnische Komponenten gegliedert, darunter \texttt{robot\_hal.hpp}, \texttt{pid\_controller.hpp} und \texttt{diff\_drive\_kinematics.hpp}.

Die Regelaufgabe läuft auf Core 1 als \texttt{controlTask} mit einer Zyklusfrequenz von

$$f_{\mathrm{ctrl}} = 50\,\mathrm{Hz}$$

Die Encoder-Impulse werden per Interrupt erfasst. Der Regler nutzt die Parametrierung

$$K_p = 0{,}4 \qquad K_i = 0{,}1 \qquad K_d = 0{,}0$$

Zur Glättung der Drehzahlmessung verwendet der Knoten einen Exponential-Moving-Average-Filter mit

$$\alpha = 0{,}3$$

Core 0 betreibt den \texttt{rclc\_executor}, publiziert Odometrie auf \texttt{/odom} mit

$$f_{\mathrm{odom}} = 20\,\mathrm{Hz}$$

und empfängt Geschwindigkeitskommandos über \texttt{/cmd\_vel}. Eine gemeinsam genutzte Datenstruktur wird durch einen Mutex geschützt, damit Regelkern und Kommunikationskern konsistente Zustände verarbeiten.

\subsection{Sensor-Knoten für I2C-Peripherie und sicherheitsnahe Signale}

Der Sensor-Knoten realisiert die Sensor- und Sicherheitsbasis. Er liest IMU, Batterieüberwachung und Kanten-Sensorik aus und stellt die Messgrößen dem ROS-2-System über micro-ROS bereit. Die Auslagerung dieser Aufgaben verhindert, dass langsame oder blockierende Sensorzugriffe die Fahrregelung beeinträchtigen.

Die I2C-Abfragen laufen auf Core 1. Da die verwendete Wire-Bibliothek nicht thread-sicher arbeitet, serialisiert ein \texttt{i2c\_mutex} alle Buszugriffe. Der Timeout beträgt

$$t_{\mathrm{mutex}} = 5\,\mathrm{ms}$$

Für die Schätzung der Gierrichtung kombiniert ein Complementary-Filter Gyroskop- und Encoderinformation mit der Gewichtung

$$98\,\%\ \text{Gyroskop} \qquad \text{und} \qquad 2\,\%\ \text{Encoder}$$

Core 0 publiziert IMU-Daten auf \texttt{/imu} mit einer Sollfrequenz von

$$f_{\mathrm{imu,soll}} = 50\,\mathrm{Hz}$$

Die tatsächlich erreichte Frequenz beträgt aufgrund von I2C-Bus-Contention etwa $33\,\mathrm{Hz}$.

Batteriedaten auf \texttt{/battery} mit

$$f_{\mathrm{battery}} = 2\,\mathrm{Hz}$$

und den Kantenstatus auf \texttt{/cliff} mit

$$f_{\mathrm{cliff}} = 20\,\mathrm{Hz}$$

Aktorkommandos, etwa für Servos, werden nicht direkt im Kommunikations-Callback ausgeführt. Stattdessen nutzt der Knoten ein Deferred-Pattern: Der Callback schreibt nur den Sollzustand in den Arbeitsspeicher, die eigentliche I2C-Ausgabe erfolgt in der regulären Ablaufsteuerung. Dieses Muster reduziert Jitter und vermeidet blockierende Buszugriffe im Callback-Kontext.

Der Ultraschallsensor HC-SR04 nutzt eine ISR-basierte, nicht blockierende Messung. Eine Interrupt-Service-Routine auf dem Echo-Pin erfasst die Laufzeit über GPIO-Register-Zugriffe. Das Trigger-Echo-Pattern ersetzt das blockierende \texttt{pulseIn()} und gibt die Rechenzeit des Sensor-Tasks für IMU-Abfragen frei. Die Entfernungsdaten erscheinen auf \texttt{/range/front} mit einer Rate von $10\,\mathrm{Hz}$.

\section{Hostseitige ROS-2-Integration}

Die hostseitige Software läuft auf dem Raspberry Pi 5 mit ROS 2 Humble in einer Container-Umgebung. Diese Ebene bündelt Lokalisierung und Kartierung, Navigation, Benutzeroberfläche, Sicherheitslogik und vorbereitete Interaktionsfunktionen. Die Implementierung hält damit die Trennung zwischen Low-Level-Regelung und rechenintensiver Verarbeitung aus dem Systementwurf ein.

\subsection{micro-ROS-Anbindung der Low-Level-Knoten}

Die beiden ESP32-S3-Knoten sind nicht als proprietäre Peripherie gekoppelt, sondern als micro-ROS-Teilnehmer in den ROS-2-Verbund eingebunden. Zwei getrennte \texttt{micro\_ros\_agent}-Prozesse binden den Drive-Knoten über \texttt{/dev/amr\_drive} und den Sensor-Knoten über \texttt{/dev/amr\_sensor} mit jeweils $921600\,\mathrm{Bd}$ an. Die serielle Kopplung folgt dem in Kapitel 3 begründeten Ziel, den Regelpfad deterministisch und störarm zu halten.

Zusätzlich synchronisiert \texttt{rmw\_uros\_sync\_session(1000)} die Zeitbasis der Mikrocontroller mit der Host-Uhr. Diese Synchronisation ist für konsistente Zeitstempel von Odometrie, Sensorik und Transformationsketten erforderlich.

\subsection{CAN-Bus als redundanter Kommunikationspfad}

Zusätzlich zur UART-basierten micro-ROS-Anbindung übertragen beide ESP32-S3-Knoten ihre Daten parallel über einen CAN-Bus mit $1\,\mathrm{Mbit/s}$ (ISO 11898, SN65HVD230-Transceiver). Die CAN-Sends laufen auf Core 1, damit sie unabhängig vom micro-ROS Agent arbeiten. Der Sensor-Knoten sendet Cliff-Signale (CAN-ID 0x120, 20 Hz) und Unterspannungsereignisse (CAN-ID 0x141) direkt an den Drive-Knoten. Der Drive-Knoten empfängt diese Frames im Regelungs-Task und setzt bei Schutzauslösung die Sollgeschwindigkeiten auf null. Dieser Pfad arbeitet unabhängig vom Raspberry Pi und bildet damit einen hardwarenahen Sicherheitspfad. Auf dem Raspberry Pi empfängt ein optionaler \texttt{can\_bridge\_node} die CAN-Frames via SocketCAN und republiziert sie als ROS-2-Topics.

\subsection{Integration von Lokalisierung und Kartierung sowie Navigation}

Die Implementierung bindet LiDAR, Odometrie und IMU in die hostseitige Verarbeitung ein. Auf dieser Grundlage laufen die Knoten für Lokalisierung und Kartierung sowie Navigation. Der Host verarbeitet damit genau die Funktionsgruppen, die in der Roadmap den Ebenen „Lokalisierung und Kartierung" sowie „Navigation" zugeordnet sind. Die Low-Level-Knoten liefern dafür Zustandsgrößen und sicherheitsnahe Signale, übernehmen jedoch keine globale Zielplanung.

Die Aufgabentrennung hat eine klare Folge: Der Host darf Bewegungsziele, Kartenbezug und Pfadverfolgung berechnen, der Fahrkern setzt ausschließlich freigegebene Sollgrößen um. Damit bleibt die Umsetzungslogik auch bei höherer Systemlast stabil.

\section{Sicherheitslogik und Kommandopfad}

Die Implementierung enthält eine eigene Sicherheitslogik, die reguläre Fahrkommandos überstimmen kann. Diese Logik setzt die in der Roadmap geforderte Trennung zwischen Navigation, Bedienung und Schutzreaktion technisch um.

\subsection{Cliff-Sicherheitsmultiplexer}

Der \texttt{cliff\_safety\_node} bildet den sicherheitsnahen Multiplexer für Kanten- und Hinderniserkennung. Er abonniert \texttt{/cliff} sowie \texttt{/range/front} und überwacht gleichzeitig die eingehenden Fahrkommandos. Erkennt der Sensor-Knoten eine Kante oder unterschreitet die Ultraschall-Distanz $80\,\mathrm{mm}$, unterbricht der Multiplexer den regulären Kommandopfad und sendet stattdessen einen Stopp auf das finale Topic \texttt{/cmd\_vel}. Die Freigabe erfolgt erst bei einer Distanz über $120\,\mathrm{mm}$ (Hysterese).

Der sichere Halt entspricht im Auslösefall der Vorgabe

$$v = 0\,\mathrm{m/s}$$

Diese Umsetzung verschiebt die Kantenerkennung bewusst vor reguläre Navigationsentscheidungen. Eine erkannte Kante ist kein normaler Navigationsfehler, sondern ein Schutzereignis mit Vorrang.

\subsection{Freigegebene Kommandokette}

Die Implementierung folgt nicht dem Muster direkter Motoransteuerung durch beliebige Eingaben. Stattdessen entsteht der Fahrbefehl aus einer freigegebenen Kommandokette. Autonome Navigation, manuelle Bedienung und spätere Sprachbefehle bleiben logisch getrennte Eingangsquellen. Erst nach Freigabe darf ein Ziel- oder Moduskommando in einen Fahrbefehl überführt werden.

Die zulässige Kette lautet damit konzeptionell:

$$\text{Interaktion} \rightarrow \text{Freigabelogik} \rightarrow \text{Missionskommando} \rightarrow \text{Navigation oder Bedienung} \rightarrow /\text{cmd\_vel} \rightarrow \text{Fahrkern}.$$

Diese Struktur verhindert, dass eine spätere Sprachschnittstelle die Sicherheitslogik umgeht. Gleichzeitig bleibt die Bedien- und Leitstandsebene handlungsfähig, weil sie den Betriebszustand beobachten und freigegebene Kommandos auslösen kann.

\section{Bedien- und Leitstandsebene}

Die Bedien- und Leitstandsebene wurde als eigenständiger Implementierungsblock aufgebaut. Sie stellt Telemetrie, Videostream, manuelle Bedienung und Zustandsanzeige bereit und dient damit sowohl dem Betrieb als auch der Diagnose. Diese Ebene greift nicht direkt in die Motorregelung ein, sondern nutzt die definierte Kommandokette.

\subsection{Dashboard-Brücke und Kommunikationspfade}

Das System stellt eine browserbasierte Benutzeroberfläche auf Basis von Vite und React bereit. Ein maßgeschneiderter \texttt{dashboard\_bridge} verbindet ROS-2-Datenpfade mit der Benutzeroberfläche. Im Unterschied zu einer generischen \texttt{rosbridge\_suite} reduziert diese Brücke den Protokoll-Overhead für den konkreten Anwendungsfall.

Die Brücke kombiniert einen WebSocket-Server für Telemetrie auf Port

$$9090$$

mit einem HTTP-Endpunkt für den MJPEG-Stream auf Port

$$8082$$

Die Benutzeroberfläche ist auf dem Host über Port

$$5173$$

erreichbar. Diese Aufteilung trennt zustandsarme Telemetrie von bandbreitenintensivem Videostream.

\subsection{Betriebsfunktionen der Benutzeroberfläche}

Die Benutzeroberfläche visualisiert LiDAR-Daten, Kamerabild, Telemetrie und manuelle Bedienelemente. Damit bildet sie die in der Roadmap geforderte Bedien- und Leitstandsebene mit Diagnosefunktion ab. Die manuelle Steuerung erfolgt über einen Joystick. Zusätzlich stehen Zustandsanzeigen zur Verfügung, die den aktuellen Betriebszustand nachvollziehbar machen.

Ein dreistufiger Deadman-Mechanismus sichert die Teleoperation ab. Die erste Stufe liegt im Browser und sendet beim Loslassen des Joysticks sofort einen Stopp. Die zweite Stufe liegt in der Dashboard-Brücke und stoppt die Ausgabe bei ausbleibendem WebSocket-Heartbeat nach mehr als

$$300\,\mathrm{ms}$$

Die dritte Stufe liegt in der Firmware und greift bei ausbleibenden \texttt{cmd\_vel}-Nachrichten nach mehr als

$$500\,\mathrm{ms}$$

ein. Die drei Stufen wirken nacheinander und erhöhen damit die Robustheit der manuellen Bedienung.

\section{Erweiterungen für Vision, Audio und Sprachschnittstelle}

Die Ebene der intelligenten Interaktion wurde nur teilweise als Kernfunktion umgesetzt. Realisiert wurden die Kameraanbindung, der hardwarebeschleunigte Vision-Pfad und ein Audiopfad. Die Sprachschnittstelle bleibt als vorbereitete Erweiterung anschlussfähig, ist jedoch nicht Teil des fahrkritischen Kernpfads.

\subsection{Hybride Vision-Pipeline}

Die Vision-Verarbeitung wurde aus dem ROS-Echtzeitgraphen ausgelagert, damit Kartierung, Navigation und Sicherheitslogik nicht durch Bildverarbeitung blockiert werden. Der \texttt{v4l2\_camera\_node} erfasst den Kamerastrom. Die Dashboard-Brücke stellt daraus einen MJPEG-Stream bereit. Ein nativer Host-Prozess \texttt{host\_hailo\_runner.py} führt die Objekterkennung auf dem Hailo-8L aus und überträgt erkannte Begrenzungsrahmen per UDP auf Port

$$5005$$

zurück in die Container-Umgebung.

Zusätzlich ist eine asynchrone Semantikstufe mit \texttt{gemini\_semantic\_node} vorgesehen. Diese Stufe bleibt bewusst vom Fahrkern entkoppelt und ergänzt den Wahrnehmungspfad nur oberhalb der Sicherheits- und Freigabelogik.

\subsection{Audioausgabe und vorbereitete Sprachintegration}

Der Audiopfad unterstützt definierte Rückmeldungen an die Bedien- und Leitstandsebene. Die Roadmap ordnet zusätzlich eine Sprachschnittstelle mit ReSpeaker Mic Array v2.0 der Ebene der intelligenten Interaktion zu. Für die Projektarbeit ist entscheidend, dass diese Erweiterung architektonisch vorbereitet bleibt: Sprachbefehle dürfen ausschließlich über Intent-Erkennung, Freigabelogik und Missionskommando auf das System wirken.

Damit gilt auch für die Erweiterung der Grundsatz, dass Sprachverarbeitung keine direkte Motoransteuerung auslösen darf. Die spätere Integrationskette lautet daher:

$$\text{Sprachbefehl} \rightarrow \text{Intent} \rightarrow \text{Freigabelogik} \rightarrow \text{Missionskommando} \rightarrow \text{Navigation, Leitstand oder Audio}.$$

\section{Ergebnis der Implementierung}

Die Implementierung überführt den Entwurf in eine funktionsfähige, klar getrennte Systemstruktur. Zwei ESP32-S3-Knoten realisieren Fahrkern sowie Sensor- und Sicherheitsbasis. Der Raspberry Pi 5 bündelt ROS 2, Lokalisierung und Kartierung, Navigation, Sicherheitslogik und die Bedien- und Leitstandsebene. Kamera, Hailo-8L und Audiopfad erweitern das System oberhalb des fahrkritischen Kerns. Damit liegt eine Implementierung vor, die die Roadmap-Themen in eine konsistente technische Gesamtstruktur überführt und die spätere Validierung gezielt vorbereitet.

\chapter{Validierung und Bewertung der Ergebnisse}

\section{Leitfrage, Prüfumfang und Datengrundlage}

Die Leitfrage dieses Kapitels lautet: In welchem Maß erfüllt das implementierte System die Kernanforderungen an Fahrkern, Sensor- und Sicherheitsbasis, Lokalisierung und Kartierung, Navigation sowie Sicherheitslogik unter den Randbedingungen eines Innenraums?

Die Validierung stützt sich auf projektinterne Messdaten aus einem $10\,\mathrm{m} \times 10\,\mathrm{m}$ großen Testareal mit statischen und dynamischen Hindernissen. Die Referenzpositionen wurden mit einem Laserdistanzmessgerät erfasst. Für die Offline-Analyse zeichnete \texttt{rosbag2} die Topics \texttt{/odom}, \texttt{/scan}, \texttt{/cmd\_vel} und \texttt{/tf} auf. Die Datengrundlage verbindet damit Bewegungsdaten des Fahrkerns mit Zustandsdaten aus Lokalisierung und Kartierung sowie der Navigation.

Der Prüfumfang gliedert sich in drei Ebenen. Die erste Ebene verifiziert den Fahrkern sowie die Sensor- und Sicherheitsbasis. Die zweite Ebene validiert Lokalisierung und Kartierung, Navigation und Sicherheitslogik im Gesamtsystem. Die dritte Ebene betrachtet erweiterte Funktionen wie Vision und Docking. Diese dritte Ebene gehört nicht zum zwingenden Kern der Roadmap, ist für die technische Einordnung des Prototyps jedoch relevant.

Tabelle 6.1 ordnet die Messgrößen den Kernzielen der Arbeit zu.

\begin{table}[htbp]
  \centering
  \caption{Zuordnung der Messgrößen zu den Kernzielen (Tabelle 6.1)}
  \small
  \begin{tabularx}{\textwidth}{l l L}
    \toprule
    \textbf{Prüfbereich} & \textbf{Messgröße} & \textbf{Kriterium} \\
    \midrule
    Fahrkern & Regeltakt, Jitter, Odometriefehler & reproduzierbare Grundbewegung und deterministische Verarbeitung \\
    \addlinespace
    Sensor- und Sicherheitsbasis & Kantenreaktionszeit, Sensorkette & sicherer Halt bei erkannter Gefahr \\
    \addlinespace
    Lokalisierung und Kartierung & Absolute Trajectory Error, Kartenqualität & belastbare Karte und Re-Lokalisierungsfähigkeit \\
    \addlinespace
    Navigation & Zielabweichung, Recovery-Verhalten, Passierabstand & sichere Zielanfahrt auf Kartenbasis \\
    \addlinespace
    Sicherheitslogik & Übersteuerung konkurrierender Kommandos & Vorrang des sicheren Zustands \\
    \addlinespace
    Erweiterungsfunktionen & Inferenzzeit, Docking-Toleranz, Erfolgsquote & Einordnung von Vision und Docking außerhalb des Kernpfads \\
    \bottomrule
  \end{tabularx}
\end{table}

\section{Verifikation von Fahrkern sowie Sensor- und Sicherheitsbasis}

\subsection{Deterministische Verarbeitung des Fahrkerns}

Der Fahrkern muss Bewegungsbefehle mit konstantem Zeitverhalten umsetzen. Gemessen wurde deshalb der Regelzyklus der PID-Schleife auf dem Drive-Knoten. Die Schleife hielt einen Takt von $50\,\mathrm{Hz}$ bei einem Jitter von unter $2\,\mathrm{ms}$ ein. Der Datenverlust der micro-ROS-Kommunikation lag unter $0{,}1\,\%$.

Das Ergebnis belegt eine hinreichend deterministische Verarbeitung für den Fahrkern. Die Trennung zwischen Drive-Knoten und Sensor-Knoten erfüllt damit die zentrale Architekturabsicht, blockierende Sensorzugriffe von der Motorregelung fernzuhalten. Das Ergebnis stützt insbesondere die nichtfunktionale Anforderung N01.

\subsection{Odometrie und Kalibrierung}

Die Grundbewegung des Fahrkerns wurde über das UMBmark-Verfahren und eine distanzbasierte Referenzfahrt geprüft. Zehn bidirektionale Testläufe zeigten vor der Kalibrierung deutliche Drift-Cluster. Nach der Kalibrierung mit einem effektiven Raddurchmesser von $65{,}67\,\mathrm{mm}$ und angepasster Spurbreite sank der systematische Fehler um den Faktor $10$ bis $20$.

Für eine Geradeausfahrt über $1{,}0\,\mathrm{m}$ reduzierte sich der laterale Fehler bei aktivierter IMU-Heading-Korrektur von $3{,}6\,\mathrm{cm}$ auf $2{,}1\,\mathrm{cm}$. Ohne IMU-Korrektur betrug der Gierfehler $+5{,}57^\circ$ (Gyro) bzw.\ $+2{,}91^\circ$ (Odometrie) und überschritt damit das Akzeptanzkriterium von $5^\circ$ knapp. Mit aktivierter IMU-Korrektur sank der Gyro-Gierfehler auf $+0{,}06^\circ$ bei einem Odometrie-Gierfehler von $+1{,}49^\circ$. Der Odometrie-Distanzfehler lag bei $0{,}5\,\%$ ohne bzw.\ $0{,}3\,\%$ mit IMU-Korrektur.

In einem ergänzenden Rotationstest sollte der Fahrkern eine Drehung um $360^\circ$ ausführen. Die tatsächlich erreichte Rotation betrug $358{,}1^\circ$, der Restfehler somit $1{,}88^\circ$. Der gemessene Gyro-Bias lag bei $-0{,}012\,\mathrm{deg/s}$ über eine Testdauer von $14{,}4\,\mathrm{s}$. Das Ergebnis liegt deutlich unter dem Akzeptanzkriterium von $5^\circ$ und belegt die Reproduzierbarkeit der Rotation nach Kalibrierung.

Das Ergebnis zeigt: Der Fahrkern erreicht nach Kalibrierung sowohl für Geradeausfahrt als auch für Rotation eine reproduzierbare Grundbewegung. Ohne IMU-Korrektur verfehlt die Geradeausfahrt das Heading-Kriterium knapp, mit IMU-Korrektur wird es klar eingehalten. Für die Zielgröße einer belastbaren Innenraum-Navigation ist die Kalibrierung einschließlich IMU-Heading-Korrektur damit nicht optional, sondern Voraussetzung.

\subsection{Reaktionszeit der Sicherheitskette}

Die Sicherheitskette wurde über manuelle Kantenüberfahrten bei einer Fahrgeschwindigkeit von $0{,}2\,\mathrm{m/s}$ geprüft. Gemessen wurde die vollständige End-to-End-Latenz von der Detektion am Kanten-Sensor über das Topic \texttt{/cliff}, die Verarbeitung im \texttt{cliff\_safety\_node} und den daraus resultierenden Stoppbefehl an den Drive-Knoten. Die gesamte Reaktionszeit betrug $2{,}0\,\mathrm{ms}$ bei einem Bremsweg von $1{,}0\,\mathrm{cm}$.

Das Ergebnis erfüllt die zentrale Forderung der Sensor- und Sicherheitsbasis: Ein erkannter Gefahrenzustand führt reproduzierbar in einen sicheren Halt. Die End-to-End-Latenz von $2{,}0\,\mathrm{ms}$ liegt deutlich unter einem regulären Replanning-Zyklus der Navigation. Für Kantenereignisse greift damit der Schutzpfad vor dem Komfortpfad.

\subsection{Einzelbewertung der Sensorik}

Die Sensor- und Sicherheitsbasis umfasst neben dem Kanten-Sensor weitere Subsysteme, deren Leistungsfähigkeit einzeln bewertet wurde. Die Messungen entstanden unter kontrollierten Laborbedingungen mit definierten Referenzabständen und -winkeln.

\paragraph{Ultraschall (HC-SR04).}
Der Ultraschallsensor erreichte eine Messrate von $9{,}2\,\mathrm{Hz}$ bei einem nutzbaren Bereich von $0{,}02$ bis $4{,}0\,\mathrm{m}$ und einem Öffnungswinkel von $0{,}26\,\mathrm{rad}$. Bei einem Referenzabstand von $23\,\mathrm{cm}$ betrug die Genauigkeit $0{,}8\,\%$ mit einer Wiederholbarkeit von $\sigma = 1{,}6\,\mathrm{mm}$. Die ISR-basierte Auswertung vermeidet blockierende Wartezeiten in der Sensor-Task und belastet die IMU-Abtastung nicht.

\paragraph{Cliff-Sensor (MH-B IR).}
Der Cliff-Sensor erreichte eine Rate von $16{,}8\,\mathrm{Hz}$ bei null Fehlalarmen über die gesamte Testreihe. Die Erkennungszeit für einen Kantenübergang betrug $1\,\mathrm{ms}$. Das Ergebnis bestätigt die Eignung als sicherheitskritischer Eingang für die Cliff-Safety-Kette.

\paragraph{IMU (MPU-6050).}
Die IMU lieferte eine Datenrate von $30$ bis $35\,\mathrm{Hz}$. Die Gyro-Drift betrug $0{,}463\,\mathrm{deg/min}$, der Beschleunigungs-Bias $0{,}43\,\mathrm{m/s^2}$. In einem $90^\circ$-Rotationstest erreichte der Fahrkern $88{,}5^\circ$, was einem Fehler von $1{,}45^\circ$ entspricht. Der gemessene Gyro-Drift bleibt für die Heading-Korrektur im Innenraumbetrieb tragbar, da die SLAM-Korrektur langfristige Drift kompensiert.

\paragraph{Batterie (INA260).}
Die Batterieüberwachung erfasst Spannung, Strom und Leistung mit $2\,\mathrm{Hz}$ über den I2C-Bus. Die Unterspannungsabschaltung bei $9{,}5\,\mathrm{V}$ mit einer Hysterese von $0{,}3\,\mathrm{V}$ wurde im Betrieb verifiziert und löst einen CAN-basierten Shutdown aus.

Zusammen mit der Cliff-Safety-Kette belegen diese Einzelmessungen, dass die Sensor- und Sicherheitsbasis alle sicherheitskritischen Eingänge mit ausreichender Rate und Genauigkeit bereitstellt.

\section{Validierung von Lokalisierung, Kartierung und Navigation}

\subsection{Lokalisierung und Kartierung im Innenraum}

Lokalisierung und Kartierung wurden über den Absolute Trajectory Error, kurz ATE, und die Wiederholbarkeit der erzeugten Karte bewertet. Der mittlere ATE der \texttt{slam\_toolbox} betrug $0{,}16\,\mathrm{m}$. Das Akzeptanzkriterium lag bei $0{,}20\,\mathrm{m}$.

Der gemessene ATE unterschreitet das Akzeptanzkriterium um $0{,}04\,\mathrm{m}$. Für einen kostengünstigen Innenraum-Prototyp ist das Ergebnis ausreichend, um Ziele auf Kartenbasis anzufahren und nach einem Neustart wieder zu lokalisieren. Die Kernanforderung an Lokalisierung und Kartierung gilt damit auf Basis der vorliegenden Messdaten als erfüllt.

\subsection{Zielanfahrt und Bahnverfolgung}

Die Qualität der Zielanfahrt wurde über Punkt-zu-Punkt-Fahrten bewertet. Die mittlere Positionsabweichung betrug $6{,}4\,\mathrm{cm}$, die mittlere Winkelabweichung $4{,}2^\circ$.

Diese Werte zeigen eine für die Innenraumanwendung brauchbare Zielgenauigkeit. Die Abweichungen liegen in einer Größenordnung, die eine zuverlässige Navigation zwischen definierten Zielpunkten erlaubt. Die Zielanfahrt bestätigt zugleich, dass die zuvor kalibrierte Odometrie, die IMU-Korrektur und die globale Kartenreferenz im Betrieb zusammenwirken.

\subsection{Hindernisvermeidung und Recovery-Verhalten}

Die Navigation musste nicht nur Zielpunkte erreichen, sondern auf Störungen geordnet reagieren. Der Nav2-Stack umfuhr statische Hindernisse mit einem minimalen Passierabstand von $12\,\mathrm{cm}$. In Situationen ohne gültigen Pfad lösten die Recovery-Verfahren \texttt{Spin}, \texttt{BackUp} und \texttt{Wait} insgesamt $80\,\%$ der Blockaden eigenständig auf.

Das Ergebnis zeigt zwei Punkte. Erstens arbeitet die lokale Bahnverfolgung auch in Engstellen noch kontrolliert. Zweitens beseitigt das Recovery-Verhalten einen großen Teil typischer Sackgassen, jedoch nicht alle. Die Navigation erreicht damit ein belastbares, aber noch nicht vollständiges Robustheitsniveau.

\subsection{Vorrang des sicheren Zustands}

Die Sicherheitslogik wurde in einer gezielt erzeugten Konfliktsituation geprüft. Die Navigation plante einen formal gültigen Pfad über eine geöffnete Bodenluke. Sobald der Kanten-Sensor am Sensor-Knoten auslöste, überstimmte der \texttt{cliff\_safety\_node} die anliegenden Navigationskommandos. Das Fahrzeug kam $3\,\mathrm{cm}$ vor der Absturzkante zum Stillstand.

Die Messung belegt den Vorrang des sicheren Zustands vor einer weiterhin gültigen Bewegungsplanung. Für die Architektur ist dieses Ergebnis zentral: Nicht das zuletzt berechnete Bewegungsziel entscheidet über die Aktorik, sondern die Sicherheitslogik mit höherer Priorität. Damit bestätigt der Versuch die vorgesehene Hierarchie aus Navigation, Sicherheitslogik und Fahrkern.

\section{Einordnung der Bedien- und Leitstandsebene sowie der Systemlast}

Die Bedien- und Leitstandsebene bildet das Betriebswerkzeug des Systems. Für die Validierung relevant sind Beobachtbarkeit, Zustandsdarstellung und die Trennung zwischen fahrkritischem Kern und benutzerseitiger Interaktion. Die Datengrundlage dieses Kapitels enthält dafür keine eigenständige Usability-Messreihe, wohl aber Last- und Betriebsdaten.

Der Raspberry Pi 5 blieb im regulären Betrieb unter $80\,\%$ CPU-Last. Der Docker-Container für ROS-2-Integration und MJPEG-Stream lag im Mittel bei etwa $35\,\%$ Auslastung. Der hardwarebeschleunigte Vision-Prozess \texttt{host\_hailo\_runner.py} lief nativ auf dem Host und band den PCIe-Beschleuniger ohne erhebliche Zusatzlast für die CPU ein. Auf Mikrocontroller-Ebene beanspruchte die Firmware des Drive-Knotens auf Core 1 weniger als $40\,\%$ der verfügbaren Rechenkapazität. Der Regulated Pure Pursuit Controller auf dem Raspberry Pi erreichte Rechenraten von mehr als $2{.}000\,\mathrm{Hz}$ bei einem geforderten Minimum von $100\,\mathrm{Hz}$.

Diese Messdaten belegen keine Benutzerqualität im engeren Sinn, sie belegen jedoch die technische Tragfähigkeit der Ebenentrennung. Telemetrie, Video und erweiterte Vision belasten den Fahrkern nicht in kritischem Maß. Für die Bedien- und Leitstandsebene bedeutet das: Die Beobachtung des Systems bleibt möglich, ohne die fahrkritische Kette strukturell zu destabilisieren.

\section{Explorative Bewertung von Vision und Docking}

\subsection{Vision-Pipeline}

Die Vision-Pipeline gehört zur erweiterten Interaktions- und Wahrnehmungsebene oberhalb des Fahrkerns. Die gemessene Inferenzzeit des Hailo-8L mit dem Modell \texttt{yolov8s} betrug im Mittel rund $34\,\mathrm{ms}$ pro Bild. Das Projektkriterium lag bei weniger als $50\,\mathrm{ms}$ pro Bild.

Damit erfüllt die lokale Bildverarbeitung das vorgegebene Zeitkriterium. Das Ergebnis rechtfertigt die Einordnung der Vision-Pipeline als echtzeitnahe Erweiterung. Gleichwohl bleibt die Funktion außerhalb des zwingenden Sicherheits- und Navigationskerns, da aus der Objekterkennung keine direkte, unüberwachte Motoransteuerung folgt.

\subsection{Docking als Erweiterungsfunktion}

Das Docking-Verfahren verbindet ArUco-Marker-Erkennung mit der Vision-Pipeline. Nach Optimierung der Dreifach-Bedingung (Ultraschall $\leq 0{,}30\,\mathrm{m}$, Marker sichtbar, lateraler Versatz $\leq 5\,\mathrm{cm}$) erreichten alle $10$ Docking-Versuche die vorgegebene Toleranz. Der mittlere laterale Versatz betrug $0{,}73\,\mathrm{cm}$, der mittlere Orientierungsfehler $14{,}1^\circ$.

Die verbesserte Erfolgsquote von $100\,\%$ gegenüber der initialen Messung ($80\,\%$) resultiert aus der Verkürzung der Docking-Distanz von $0{,}60\,\mathrm{m}$ auf $0{,}30\,\mathrm{m}$ und der Einführung einer Fehlausrichtungs-Recovery mit Rückwärtsfahrt und Neuausrichtung. Die höhere Orientierungsstreuung ($14{,}1^\circ$ gegenüber $2{,}8^\circ$) deutet darauf hin, dass die engere Distanz den lateralen Versatz auf Kosten der Winkelgenauigkeit verbessert. Als Erweiterungsfunktion ist das Ergebnis aussagekräftig, weil es die Grenzen monokularer Bildauswertung unter variabler Beleuchtung offenlegt.

\section{Soll-Ist-Vergleich der Anforderungen}

Tabelle 6.2 fasst die Anforderungen aus Kapitel 3 mit dem Stand der vorliegenden Messdaten zusammen. Die Bewertung unterscheidet bewusst zwischen „erfüllt", „teilweise erfüllt" und „nicht geprüft". Diese Unterscheidung trennt belegte Ergebnisse von plausiblen, aber noch unvollständig vermessenen Annahmen.

\begin{table}[htbp]
  \centering
  \caption{Soll-Ist-Vergleich der Anforderungen (Tabelle 6.2)}
  \small
  \begin{tabularx}{\textwidth}{l L l L}
    \toprule
    \textbf{ID} & \textbf{Anforderung} & \textbf{Bewertung} & \textbf{Begründung} \\
    \midrule
    F01 & Fahrkern mit reproduzierbarer Grundbewegung & erfüllt & Geradeausfahrt und $360^\circ$-Rotation sind nach Kalibrierung mit IMU-Korrektur belegt. Lateraler Fehler $2{,}1\,\mathrm{cm}$, Heading-Fehler $0{,}06^\circ$, Rotationsfehler $1{,}88^\circ$. \\
    \addlinespace
    F02 & Sensor- und Sicherheitsbasis mit priorisierten Schutzfunktionen & erfüllt & End-to-End-Latenz der Cliff-Safety-Kette $2{,}0\,\mathrm{ms}$, Bremsweg $1{,}0\,\mathrm{cm}$. Ultraschall ($9{,}2\,\mathrm{Hz}$, $0{,}8\,\%$ Fehler), Cliff ($16{,}8\,\mathrm{Hz}$, $0$ Fehlalarme), IMU ($30$--$35\,\mathrm{Hz}$, $1{,}45^\circ$ Rotationsfehler) und Batterie einzeln bewertet. \\
    \addlinespace
    F03 & Lokalisierung und Kartierung für den Innenraum & erfüllt & Der mittlere ATE von $0{,}16\,\mathrm{m}$ unterschreitet das Kriterium von $0{,}20\,\mathrm{m}$. \\
    \addlinespace
    F04 & Navigation mit Missionslogik und Recovery-Verhalten & teilweise erfüllt & Zielanfahrt, Hindernisvermeidung und Recovery-Verhalten sind belegt. Die Recovery-Verfahren lösen $80\,\%$ der Blockaden, aber nicht alle. \\
    \addlinespace
    F05 & Bedien- und Leitstandsebene als Betriebswerkzeug & teilweise erfüllt & Lastdaten und Betriebsfähigkeit stützen die technische Umsetzbarkeit. Eine eigenständige Messreihe zur Bedienqualität wurde nicht durchgeführt. \\
    \addlinespace
    F06 & Freigabelogik mit sicheren Zustandsübergängen & teilweise erfüllt & Der Vorrang des sicheren Zustands ist über den Cliff-Versuch belegt. Eine vollständige Prüfung aller Kommandoklassen liegt nicht vor. \\
    \addlinespace
    F07 & Sprachschnittstelle als anschlussfähige Erweiterung & nicht geprüft & Die Sprachschnittstelle ist architektonisch vorbereitet, aber nicht Gegenstand der Kernvalidierung. \\
    \addlinespace
    N01 & Deterministische Verarbeitung im Fahrkern & erfüllt & $50\,\mathrm{Hz}$ Regeltakt, Jitter unter $2\,\mathrm{ms}$, geringe Paketverluste. \\
    \addlinespace
    N02 & Modulare und wartbare Systemstruktur & nicht geprüft & Die Architektur ist umgesetzt, wurde in diesem Kapitel aber nicht mit einem separaten Wartbarkeitskriterium vermessen. \\
    \addlinespace
    N03 & Robuste Kommunikation und Fehlerbehandlung & teilweise erfüllt & Geringer Datenverlust ist belegt. Ein vollständiger Wiederanlauftest der Kommunikationskette wird in diesem Kapitel nicht dokumentiert. \\
    \addlinespace
    N04 & Beobachtbarkeit und Nachvollziehbarkeit & teilweise erfüllt & \texttt{rosbag2}, Telemetrie und Lastdaten zeigen Beobachtbarkeit. Ein explizites Diagnosekriterium wurde nicht separat vermessen. \\
    \addlinespace
    N05 & Erweiterbarkeit der Interaktionsschicht & nicht geprüft & Die Architektur erlaubt Vision, Audio und Sprachschnittstelle, wurde dafür aber nicht mit einem formalen Kriterium getestet. \\
    \bottomrule
  \end{tabularx}
\end{table}

\section{Schlussfolgerung aus der Validierung}

Die Validierung zeigt ein klares Gesamtbild. Der Prototyp erfüllt die Kernziele für den Fahrkern mit reproduzierbarer Grundbewegung, für die Sensor- und Sicherheitsbasis, für Lokalisierung und Kartierung, für die sichere Navigation auf Kartenbasis und für die deterministische Verarbeitung im Fahrkern. Die Sicherheitslogik greift im Konfliktfall zuverlässig vor der Navigation ein. Damit ist die zentrale Architekturentscheidung der Arbeit technisch bestätigt.

Die Sensor- und Sicherheitsbasis ist durch die Einzelbewertung von Ultraschall, Cliff-Sensor, IMU und Batterieüberwachung vollständig belegt. Die End-to-End-Latenz der Cliff-Safety-Kette von $2{,}0\,\mathrm{ms}$ und der Bremsweg von $1{,}0\,\mathrm{cm}$ bei $0{,}2\,\mathrm{m/s}$ bestätigen die Schutzfunktion unter realistischen Betriebsbedingungen. Die Bedien- und Leitstandsebene ist technisch tragfähig, aber nicht mit einer eigenständigen Messreihe zur Bedienqualität bewertet. Die Sprachschnittstelle bleibt eine vorbereitete Erweiterung und gehört folgerichtig nicht zum belegten Kernumfang.

Für die Projektarbeit folgt daraus ein belastbarer Schluss: Das System erreicht eine funktionsfähige und sicher priorisierte Innenraum-Mobilität auf Basis einer verteilten, kostengünstigen Architektur. Fahrkern und Sensor- und Sicherheitsbasis erfüllen nach Kalibrierung und IMU-Korrektur sämtliche Akzeptanzkriterien. Die größten Verbesserungshebel liegen nicht mehr im Fahrkern oder in der Sensorik, sondern in der Robustheit des Recovery-Verhaltens und in der systematischen Vermessung der erweiterten Interaktionsschicht.

\chapter{Fazit und Ausblick}

\section{Beantwortung der Leitfrage und der Projektfragen}

Die Leitfrage dieser Arbeit lautet, ob sich ein kostengünstiges autonomes mobiles Robotersystem so auslegen lässt, dass Fahrkern, Sensor- und Sicherheitsbasis, Lokalisierung und Kartierung, Navigation sowie Bedien- und Leitstandsebene im Innenraum reproduzierbar zusammenarbeiten. Die Ergebnisse aus Kapitel 6 beantworten diese Leitfrage im Kern positiv. Der Prototyp erreicht eine funktionsfähige und sicher priorisierte Innenraum-Mobilität auf Basis einer verteilten Architektur aus zwei ESP32-S3-Mikrocontrollern und einem Raspberry Pi 5. Gleichzeitig zeigt die Validierung, dass nicht alle Akzeptanzkriterien in gleicher Tiefe nachgewiesen sind und dass erweiterte Funktionen nicht mit dem gleichen Reifegrad vorliegen wie der fahrkritische Kern.

Die Arbeit zielt damit nicht auf den vollständigen Nachweis industrieller Serienreife, sondern auf den belastbaren Nachweis einer tragfähigen Systemarchitektur für einen kostengünstigen Innenraum-AMR. Diese Einordnung entspricht auch der Roadmap: Zuerst müssen Fahrkern, Sensor- und Sicherheitsbasis, Lokalisierung und Kartierung, Navigation sowie Bedien- und Leitstandsebene belastbar arbeiten. Erst danach folgen Erweiterungen wie Vision oder Sprachschnittstelle.

\subsection{PF1 — Fahrkern und verteilte Echtzeitarchitektur}

Die erste Projektfrage untersucht, ob sich ein Fahrkern auf Basis einer verteilten ESP32-S3-Architektur so realisieren lässt, dass Motorregelung, Odometrie und die Verarbeitung von Geschwindigkeitskommandos deterministisch arbeiten. Die Validierung zeigt für die PID-Regelschleife einen Regeltakt von $50\,\mathrm{Hz}$ bei einem Jitter von unter $2\,\mathrm{ms}$. Der Datenverlust der micro-ROS-Kommunikation blieb unter $0{,}1\,\%$. Diese Messwerte belegen, dass die Trennung von Drive-Knoten und Sensor-Knoten blockierende Peripheriezugriffe vom Fahrkern fernhält und damit die zentrale Architekturabsicht erfüllt.

PF1 ist damit positiv beantwortet. Eine verteilte Echtzeitarchitektur mit UART-basierter micro-ROS-Anbindung kann den Fahrkern für einen kostengünstigen Innenraum-AMR deterministisch stabil betreiben, sofern die Aufgabenteilung zwischen Motorregelung und Sensorverarbeitung physisch getrennt bleibt.

\subsection{PF2 — Sensor- und Sicherheitsbasis sowie Lokalisierung und Kartierung}

Die zweite Projektfrage untersucht, welchen Beitrag Odometrie-Kalibrierung, IMU-Integration, Kanten-Erkennung und LiDAR-basierte Kartierung zur belastbaren Zustands- und Umgebungsbasis leisten. Die Validierung zeigt drei zentrale Ergebnisse. Erstens reduzierte die UMBmark-Kalibrierung den systematischen Odometriefehler um den Faktor $10$ bis $20$. Zweitens sank der laterale Drift auf einer Referenzfahrt über $1{,}0\,\mathrm{m}$ mit IMU-Korrektur von $3{,}6\,\mathrm{cm}$ auf $2{,}1\,\mathrm{cm}$, während der Gierfehler von $5{,}57^\circ$ auf $0{,}06^\circ$ sank. Drittens erreichte die Kartierung einen mittleren Absolute Trajectory Error von $0{,}16\,\mathrm{m}$ und blieb damit unter dem Akzeptanzkriterium von $0{,}20\,\mathrm{m}$. Zusätzlich reagierte die Sicherheitskette bei einem Kantenereignis in unter $50\,\mathrm{ms}$ (gemessen: $2{,}0\,\mathrm{ms}$).

PF2 ist damit weitgehend positiv beantwortet. Odometrie-Kalibrierung, IMU-Stützung, LiDAR-basierte Kartierung und priorisierte Sicherheitslogik bilden gemeinsam eine belastbare Grundlage für den Innenraumbetrieb. Die Einzelbewertung aller Sensorpfade einschließlich Ultraschall, Batterieüberwachung und IMU-Plausibilisierung ist mit $11/11$ bestandenen Testfällen abgeschlossen (vgl.\ Kapitel~6, Abschnitt~6.2.4).

\subsection{PF3 — Navigation sowie Bedien- und Leitstandsebene}

Die dritte Projektfrage untersucht, ob sich auf dieser Basis eine Systemarchitektur aufbauen lässt, die Zielanfahrten, Recovery-Verhalten, Telemetrie, manuelle Eingriffe und Audio-Rückmeldungen zusammenführt, ohne die Freigabelogik zu unterlaufen. Die Validierung zeigt für die Navigation eine mittlere Positionsabweichung von $6{,}4\,\mathrm{cm}$ und eine mittlere Winkelabweichung von $4{,}2^\circ$. Der Nav2-Stack umfuhr statische Hindernisse mit einem minimalen Passierabstand von $12\,\mathrm{cm}$. Die Recovery-Verfahren lösten $80\,\%$ der Blockaden eigenständig. In einer gezielt erzeugten Konfliktsituation überstimmte die Sicherheitslogik die gültige Bewegungsplanung, und das System kam $3\,\mathrm{cm}$ vor einer Absturzkante zum Stillstand. Diese Messwerte belegen, dass die Navigation funktional arbeitet und dass der sichere Zustand Vorrang vor dem zuletzt berechneten Bewegungsziel hat.

Für die Bedien- und Leitstandsebene liegt der Nachweis vor allem auf technischer Ebene vor. Telemetrie, Video und Diagnosefunktionen lassen sich betreiben, ohne den Fahrkern strukturell zu destabilisieren. Eine eigenständige Messreihe zur Bedienqualität, zur Verständlichkeit von Zustandsanzeigen oder zur Wirksamkeit von Audio-Rückmeldungen wurde jedoch nicht durchgeführt. PF3 ist deshalb im Kern positiv, in Bezug auf die Benutzerqualität jedoch nur teilweise beantwortet. Die Freigabelogik ist architektonisch richtig eingeordnet, aber noch nicht über alle Kommandoklassen systematisch vermessen.

\subsection{Kernaussage der Arbeit}

Die Arbeit zeigt, dass eine kostengünstige AMR-Architektur für strukturierte Innenräume technisch tragfähig ist, wenn der Fahrkern deterministisch ausgelegt, die Sensor- und Sicherheitsbasis priorisiert behandelt und die Bedien- und Leitstandsebene klar von der fahrkritischen Kette getrennt wird. Der belastbare Erkenntnisgewinn liegt damit weniger in einer einzelnen Demonstrationsfunktion als in der hierarchischen Kopplung der Systemebenen.

\section{Kritische Würdigung und Limitationen}

\subsection{Vollständigkeit des Nachweises}

Die Validierung bestätigt den fahrkritischen Kern, aber nicht alle Anforderungen mit gleicher Tiefe. Für den Fahrkern wurde die $360^\circ$-Rotation gemessen und ergab einen Ist-Wert von $358{,}1^\circ$ bei einem Restfehler von $1{,}88^\circ$ (vgl.\ Kapitel~6). Eine systematische Messreihe zum Nachlauf nach dem Stopp steht hingegen noch aus. In der Sensor- und Sicherheitsbasis sind alle $11/11$ Sensorik-Testfälle bestanden, darunter Kanten-Erkennung, Ultraschall, Batterieüberwachung und IMU-Plausibilisierung (vgl.\ Kapitel~6, Abschnitt~6.2.4). Für die Bedien- und Leitstandsebene fehlt eine eigenständige Untersuchung der Bedienqualität. Die Sprachschnittstelle ist in der Roadmap als Erweiterung beschrieben, gehört aber nicht zum validierten Kernumfang dieser Arbeit.

\subsection{Technische Grenzen des Prototyps}

Die Kalibrierung reduziert systematische Odometriefehler deutlich, kompensiert aber keinen Schlupf auf wechselnden oder unebenen Untergründen. Das Recovery-Verhalten erhöht die Robustheit der Navigation, beseitigt jedoch nicht alle Blockadesituationen. Die Vision-Pipeline erreicht mit dem Hailo-8L eine mittlere Inferenzzeit von rund $34\,\mathrm{ms}$ pro Bild und ist damit technisch leistungsfähig, bleibt jedoch außerhalb des sicherheitskritischen Kerns. Das Docking erreichte nach Optimierung der Dreifach-Bedingung eine Erfolgsquote von $100\,\%$ ($10/10$) und zeigte damit technische Machbarkeit als Erweiterungsfunktion. Beleuchtungsartefakte und ungünstige Ankunftsposen wirken sich weiterhin deutlich auf die monokulare Bildauswertung aus.

\subsection{Methodische Grenzen}

Das Vorgehen nach VDI 2206 strukturiert die Entwicklung des mechatronischen Gesamtsystems klar. Für die Softwareentwicklung mit wiederholten Anpassungen an ROS-2-Konfiguration, Recovery-Verhalten und Knotenstruktur entsteht jedoch ein iterativer Arbeitsmodus, der im linearen Bild des V-Modells nur eingeschränkt sichtbar wird. Die Arbeit bestätigt damit die Eignung des V-Modells als Ordnungsrahmen, nicht jedoch als vollständige Beschreibung realer Softwareiteration.

\subsection{Übertragbarkeit und wirtschaftliche Einordnung}

Die Ergebnisse gelten für strukturierte Innenräume mit kontrollierbaren Randbedingungen. Eine direkte Übertragung auf industrielle Daueranwendung, Mischverkehr mit Personen oder zertifizierungspflichtige Flurförderzeuge ist nicht zulässig. Dafür fehlen unter anderem redundante Sicherheitssensorik, formale Sicherheitsnachweise und eine robuste Hardwareintegration. Wirtschaftlich zeigt der Prototyp dennoch einen relevanten Punkt: Mit Consumer-Hardware und Open-Source-Software lässt sich ein funktional überzeugender Innenraum-AMR zu deutlich geringeren Kosten aufbauen als ein industrielles Serienprodukt. Die Differenz in Preis und Reifegrad markiert zugleich die Grenze zwischen Laborprototyp und industrieller Lösung.

\section{Ausblick und Weiterentwicklung}

\subsection{Kernpfad der Roadmap vervollständigen}

Der nächste Entwicklungsschritt sollte nicht mit zusätzlichen Demonstrationsfunktionen beginnen, sondern mit der vollständigen Schließung des Kernpfads. Dazu gehören eine erweiterte Messreihe für den Fahrkern, insbesondere zum Nachlaufverhalten nach dem Stopp, sowie eine formale Prüfung aller relevanten Zustandsübergänge der Freigabelogik. Für die Navigation ist insbesondere das Recovery-Verhalten weiter zu verbessern, bis Fehlfahrten, Sackgassen und lokale Blockaden reproduzierbar beherrscht werden. Diese Reihenfolge entspricht der Roadmap-Logik, in der quantitative Stabilität vor zusätzlicher Systemkomplexität steht.

\subsection{Bedien- und Leitstandsebene systematisch ausbauen}

Die Bedien- und Leitstandsebene ist technisch vorhanden, sollte aber als Betriebswerkzeug systematischer bewertet werden. Dafür sind Messgrößen für Telemetrie-Verzögerung, Zustandstransparenz, Eingriffszeit, Audio-Rückmeldungen und Fehlerdiagnose festzulegen. Ziel ist eine Benutzeroberfläche, die nicht nur funktioniert, sondern Betriebszustände eindeutig und belastbar vermittelt. Die Trennung zwischen Bedienung und Fahrlogik bleibt dabei unverändert, weil genau diese Trennung die Stabilität des Gesamtsystems absichert.

\subsection{Sprachschnittstelle als sichere Erweiterung integrieren}

Die Roadmap beschreibt die Sprachschnittstelle ausdrücklich als Erweiterung oberhalb des fahrkritischen Kerns. Diese Einordnung ist technisch zwingend. Sprachbefehle dürfen nicht direkt in rohe Geschwindigkeitsbefehle übergehen, sondern müssen über Intent-Erkennung, Freigabelogik und Missionskommando auf Navigation, Leitstand oder Audio wirken. Für die nächste Ausbaustufe bietet sich daher eine Architektur mit ReSpeaker Mic Array v2.0, \texttt{voice\_input\_node}, \texttt{speech\_to\_text\_node}, \texttt{voice\_intent\_node}, \texttt{voice\_command\_mux} und \texttt{text\_to\_speech\_node} an. Der erste validierbare Funktionsumfang sollte auf einen kleinen Wortschatz mit priorisiertem „Stopp", Moduswechseln und wenigen freigegebenen Missionskommandos begrenzt bleiben.

\subsection{Hardware- und Sicherheitsreife erhöhen}

Für einen robusteren Dauerbetrieb sind eine eigene Leiterplatine, definierte Steckverbinder, verbesserte elektromagnetische Verträglichkeit und eine mechanisch stabilere Sensorintegration naheliegende nächste Schritte. Zusätzlich sollte die Sicherheitskette um weitere Schutzpfade wie Bumper oder überwachte Not-Halt-Funktionen ergänzt werden. Für jeden Schritt in Richtung realer Betriebsumgebung steigt die Bedeutung normativer Anforderungen, insbesondere im Hinblick auf CE-Konformität und auf sicherheitsgerichtete Sensorik.

\section{Schluss}

Die Arbeit liefert den Nachweis, dass ein kostengünstiger Innenraum-AMR auf Basis einer verteilten Open-Source-Architektur technisch tragfähig aufgebaut und in seinen Kernfunktionen experimentell belegt werden kann. Der Fahrkern arbeitet deterministisch, die Sensor- und Sicherheitsbasis priorisiert sichere Zustände, Lokalisierung und Kartierung schaffen eine belastbare Umgebungsreferenz, und die Navigation erreicht eine praxistaugliche Zielanfahrt im Innenraum. Offene Punkte liegen vor allem in der vollständigen Abdeckung aller Akzeptanzkriterien und in der Robustheit erweiterter Interaktionsfunktionen. Genau daraus folgt der weitere Entwicklungsweg: zuerst den Kernpfad vollständig schließen, danach multimodale Erweiterungen wie Vision und Sprachschnittstelle kontrolliert und freigabelogisch abgesichert integrieren.


\appendix
\listoffigures
\listoftables

\begin{thebibliography}{99}
\end{thebibliography}

\end{document}
